
\documentclass{beamer}

\usetheme{metropolis}

% NOTE: deliberately NOT using booktabs -- all tables are plain tabular + \hline
\usepackage{xcolor}
\usepackage{array}

\hypersetup{
    colorlinks=true,
    urlcolor=blue
}

\title{Product-Key Memory as a Modular Unlearning Module}
\subtitle{Spam unlearning for generative recommendation}
\date{\today}

% ---------------------------------------------------------------------------
% Macros (extend here)
% ---------------------------------------------------------------------------
\newcommand{\pierre}[1]{\textcolor{cyan}{#1}}
\newcommand{\todo}[1]{\textcolor{red}{\textbf{TODO:} #1}}
\newcommand{\done}[1]{\textcolor{green!50!black}{#1}}
\newcommand{\open}[1]{\textcolor{orange!80!black}{#1}}

\newcommand{\If}{\ensuremath{I_t}}      % target (spam) item set
\newcommand{\Df}{\ensuremath{D_f}}      % forget set
\newcommand{\Dr}{\ensuremath{D_r}}      % retain set
\newcommand{\lf}{\ensuremath{\lambda_{\text{forget}}}}
\newcommand{\ls}{\ensuremath{\lambda_{\text{sep}}}}
\newcommand{\gf}{\ensuremath{g_f}}
\newcommand{\gr}{\ensuremath{g_r}}
% per-slot variants: \gf already carries a subscript, so \gf_{,i} would be a
% double subscript -- use these instead.
\newcommand{\gfi}{\ensuremath{g_{f,i}}}
\newcommand{\gri}{\ensuremath{g_{r,i}}}

% status markers for the plan tables
\newcommand{\stdone}{\textcolor{green!50!black}{done}}
\newcommand{\strun}{\textcolor{blue}{running}}
\newcommand{\sttodo}{\textcolor{red}{todo}}
\newcommand{\stdrop}{\textcolor{gray}{dropped}}

\begin{document}

\maketitle

% ===========================================================================
% \begin{frame}{TL;DR}
% \begin{itemize}
%   \item We can remove $\sim$99\% of spam exposure at \emph{clean-level utility}
%         by replacing two decoder FFNs with a PKM and rebuilding them from
%         retain data only.
%   \item The spam is \textbf{decoder-localised} for removal, but an
%         \textbf{encoder} PKM prevents the attack from being learned in the
%         first place ($30\times$ less pickup).
%   \item \alert{The memory collapses when trained jointly}: 32 of 262{,}144 slots. Post-hoc PKMs (trained in isolation) do not collapse.
%   \item Top-$t$ slot selection is implemented but only testable on the non-collapsed (post-hoc) models.
% \end{itemize}
% \end{frame}

% ===========================================================================
\section{Setup}

\begin{frame}{Setup: model, attack}
\textbf{Model.} TIGER: small encoder--decoder trained from
scratch on semantic-ID sequences ($d_{\text{model}}=128$, 4 encoder and 4 decoder blocks, $d_{ff}=1024$, RQ semantic IDs $L=4$, codebook width 256).

\textbf{Attack.} Bandwagon: fake users click one target item among popular
fillers. Beauty, $1\%$ poisoning, $n_{\text{target}}=1$, target strategy
\texttt{mid} unless stated.

\end{frame}

\begin{frame}{Setup: metrics, reference points}
    
\textbf{Metrics.}
\begin{itemize}
  \item \textbf{SH@$k$} spam hit rate --- fraction of test users seeing \If{}
  \item \textbf{ASI@$k$}, \textbf{TPM@$k$} --- related exposure measures
  \item \textbf{N@10} nDCG (utility), \textbf{UR} $=$ N@10 / clean N@10
  \item \textbf{removal} $= 1 - \text{SH}_{\text{unlearned}}/\text{SH}_{\text{poison}}$
\end{itemize}

\textbf{Reference points (non-PKM), dataset beauty.}
\begin{center}
\begin{tabular}{|l|c|c|}
\hline
model & SH@10 & N@10 \\
\hline
clean            & 0.00000 & 0.03763 \\
poisoned         & 0.01274 & 0.03619 \\
\hline
\end{tabular}
\end{center}
\end{frame}

% ---------------------------------------------------------------------------
\begin{frame}{Why memory? (Geva et al., 2021)}
FFNs behave like \textbf{key--value memories}:
\begin{itemize}
  \item \textbf{keys} capture input patterns
  \item \textbf{values} encode what to output (a distribution over the vocab)
  \item each FFN output is a composition of values; layers refine step by step
\end{itemize}

\medskip
\textbf{Takeaway.} Forgetting by fine-tuning overwrites internal memory
$\Rightarrow$ we want a \emph{modular stabilizer} instead of more backbone edits.

% \medskip
% Our data supports this in a specific way: spam removal only works when
% we touch \emph{decoder} FFNs --- consistent with values encoding the output
% distribution, and generation being decoder-side.
\end{frame}

% ===========================================================================
\section{Design space}

\begin{frame}{Design space: the knobs}
\begin{center}
\begin{tabular}{|l|l|}
\hline
\textbf{axis} & \textbf{options} \\
\hline
When       & trained-in $\;|\;$ post-hoc (insert into trained model) \\
Mode       & \texttt{replace} FFN $\;|\;$ \texttt{add} alongside FFN \\
Where      & encoder / decoder $\times$ block $\{0,1,2,3\}$, or sets \\
Update scope & all params $\;|\;$ \texttt{pkm\_only} $\;|\;$ values only \\
Selection  & none $\;|\;$ AF $\;|\;$ AF-IHF $\;|\;$ grad $\;|\;$ grad+dot \\
$t$        & $1 \ldots 10^4$ (bounded by \emph{live} slots) \\
Optimizer  & Adam $\;|\;$ SGD $\;|\;$ AdamW \\
Init       & \texttt{normal}$\times s$ $\;|\;$ \texttt{zeros} \\
Retain data & sampled subset $\;|\;$ full retain split \\
\hline
\end{tabular}
\end{center}

\pierre{All of these are implemented as config knobs; defaults reproduce the
PKM-free behavior.}
\end{frame}

% ---------------------------------------------------------------------------
\begin{frame}{Two variants of the idea}
\textbf{A. Trained-in (``pretrained-memory'').}
PKM is part of the model during training, so it \emph{learns} spam into memory
slots; unlearning edits those slots.

\medskip
\textbf{B. Post-hoc (``correction module'').}
PKM is inserted into an already-trained model, then adapted on clean data as a
sparse corrective layer against a frozen backbone.

\medskip
If PKM is added after training, it has no target information to remove? $\to$ removal is the forgetting, fine-tuning on the retain set repairs this

\end{frame}

% ===========================================================================
\section{Experiments done}

\begin{frame}{Headline comparison (beauty, bandwagon \texttt{mid})}
\begin{center}
\begin{tabular}{|l|c|c|c|}
\hline
method & SH@10 & removal & UR \\
\hline
poison baseline                & 0.01274 & ---    & 0.962 \\
\texttt{seif}                  & $\sim$0 & $\sim$full & 0.67--0.74 \\
\texttt{unified} ac0 (full model) & 0.00411 & 67\% & 1.013 \\
\hline
trained-in PKM, \texttt{pkm\_only} & 0.00385 & 69\% & 1.004 \\
\textbf{post-hoc d01 + rebuild}    & \textbf{0.00013} & \textbf{99\%} & \textbf{0.995} \\
\quad \emph{control:} fresh FFN, no PKM & 0.00022 & 98\% & 0.995 \\
\hline
\end{tabular}
\end{center}

\medskip
\begin{itemize}
  \item Post-hoc \textbf{dominates} everything: clean-equivalent on \emph{both}
        axes (clean is SH 0.00000 / N@10 0.03763).
  \item Trained-in \texttt{pkm\_only} matches full-model unlearning with the
        \textbf{backbone frozen} --- 0.2\% utility cost vs.\ its own clean model.
\end{itemize}
\end{frame}

% ---------------------------------------------------------------------------
\begin{frame}{Placement, post-hoc: spam is decoder-localised}
Insert PKM at \emph{one} layer, rebuild on retain data:
\begin{center}
\begin{tabular}{|l|c|c|c||l|c|c|c|}
\hline
layer & SH@10 & rem. & UR & layer & SH@10 & rem. & UR \\
\hline
d0 & 0.00139 & 89\% & 0.978 & e0 & 0.01306 & $-3\%$  & 0.978 \\
d1 & 0.00031 & 98\% & 0.958 & e1 & 0.01413 & $-11\%$ & 0.978 \\
d2 & 0.00009 & 99\% & 0.904 & e2 & 0.01288 & $-1\%$  & 0.960 \\
d3 & 0.00054 & 96\% & 0.898 & e3 & 0.01400 & $-10\%$ & 0.944 \\
\hline
\end{tabular}
\end{center}

\begin{itemize}
  \item \textbf{Every encoder layer removes nothing.} All decoder layers remove
        89--99\%.
  \item Deeper decoder $\Rightarrow$ more removal, less utility.
\end{itemize}
\end{frame}

\begin{frame}{Placement, post-hoc: spam is decoder-localised (cont.)}
\begin{itemize}
  \item \texttt{d01} (both) beats every single layer: 99\% at UR 0.995
        $\Rightarrow$ optimum is the \emph{first two} decoder FFNs.
  \item \texttt{dall} (4 layers) does \emph{not} converge $\Rightarrow$
        a destruction threshold.
  \item Need comparison with FFN reinit instead of PKM replace
\end{itemize}
\end{frame}

% ---------------------------------------------------------------------------
\begin{frame}{Placement, trained-in: encoder PKM is prophylactic}
PKM present \emph{during} training (poisoned run), spam$\times$ relative to
non-PKM (SH 0.01274):
\begin{center}
\small
\begin{tabular}{|l|c|c|c||l|c|c|c|}
\hline
layer & N@10 & SH@10 & spam$\times$ & layer & N@10 & SH@10 & spam$\times$ \\
\hline
D0 & 0.03452 & 0.00277 & 0.22 & E0 & 0.03511 & 0.00072 & 0.06 \\
D1 & 0.03607 & 0.02111 & \alert{1.66} & E1 & 0.03507 & 0.00036 & \textbf{0.03} \\
D2 & 0.03508 & 0.00648 & 0.51 & E2 & 0.03596 & 0.00778 & 0.61 \\
D3 & 0.03480 & 0.01011 & 0.79 & E3 & 0.03508 & 0.00040 & \textbf{0.03} \\
\hline
\end{tabular}
\end{center}

\end{frame}

\begin{frame}{Placement, trained-in: encoder PKM is prophylactic (cont.)}
\textbf{Encoder PKM suppresses the attack $\sim 30\times$} (E1/E3) for $\sim$3\%
utility. Exactly the opposite of the post-hoc screen --- and consistent:

\begin{center}
\emph{encoder PKM = prophylactic (prevents learning spam)}\\
\emph{decoder PKM = therapeutic (enables removing spam)}
\end{center}

% \pierre{Utility spread (0.954--0.997) is inside seed noise; the spam spread
% ($55\times$) is not.}
\end{frame}

% ---------------------------------------------------------------------------
\begin{frame}{\alert{Problem: the memory collapses}}
Slot-access diagnostic, $n_{\text{keys}}=512 \Rightarrow 262{,}144$ slots/memory:
\begin{center}
\begin{tabular}{|l|c|c|c|c|}
\hline
model & slots & touched & Jaccard & forget-excl. \\
\hline
trained-in E23D01 & 262144 & \alert{32} & \alert{1.000} & \alert{0} \\
trained-in D1only & 262144 & \alert{32} & \alert{1.000} & \alert{0} \\
trained-in D2only & 262144 & \alert{32} & \alert{1.000} & \alert{0} \\
post-hoc \texttt{replace} & 262144 & 64052 & 0.386 & 3051 \\
post-hoc \texttt{add}     & 262144 & 71547 & 0.411 & 3194 \\
\hline
\end{tabular}
\end{center}

\begin{itemize}
  \item Jointly-trained PKMs use \textbf{0.012\%} of the memory; every token and
        every head picks the \emph{same} 32 slots (routing is input-independent).
  \item \textbf{Post-hoc PKMs are healthy} --- 59\% coverage, forget/retain route
        differently, $\sim$3000 forget-exclusive slots.
\end{itemize}
\end{frame}

\begin{frame}{\alert{Problem: the memory collapses} (cont.)}
\begin{itemize}
  \item \pierre{Hypothesis: trained in isolation on a frozen backbone, the query
        net \emph{must} learn routing; in joint training the model routes
        \emph{around} the memory.}
  \item Consequence: every trained-in PKM number describes a \emph{low-rank
        layer}, not a large sparse memory.
\end{itemize}
\end{frame}

% ---------------------------------------------------------------------------
\begin{frame}{Top-$t$ slot selection: criteria}
Adapt Sparse Memory Finetuning to unlearning --- score slots, update only top $t$.
\begin{center}
\begin{tabular}{|l|l|}
\hline
criterion & score \\
\hline
AF      & forget access frequency (baseline) \\
AF-IHF  & $\text{AF}(s)\cdot\log\frac{T_r+1}{\text{HF}(s)+1}$, history $=$ retain \\
grad    & $\|\gf\| - \lambda\|\gr\|$ \\
grad+dot & $\hat{g}_f - \lambda\hat{g}_r - \mu\,\widehat{\langle \gf,\gr\rangle}$ \\
\hline
\end{tabular}
\end{center}

\textbf{Why the dot term.} The update moves along $+\gf$, so the first-order
change in the \emph{retain} loss from editing slot $i$ is
$\langle \gfi, \gri \rangle$ --- \emph{not} $\|\gri\|$. It is signed:
negative means forgetting also \emph{helps} retain.

\end{frame}

\begin{frame}{Top-$t$ slot selection: criteria}
\textbf{Verified.} Two slots with identical $\|\gf\|=\|\gr\|=12$ but dot $0$
vs.\ $+144$: ranked 1007 vs.\ 1006 by $\gf-\gr$ (indistinguishable), 1007
vs.\ 1999 (last) by grad+dot.

\medskip
\pierre{Objective is additively separable over slots $\Rightarrow$ exact top-$t$
is plain \texttt{topk}, no greedy needed.}
\end{frame}

% ---------------------------------------------------------------------------
\begin{frame}{Selection: gradient signal separates, access counts do not}
Mean per-slot forget/retain gradient cosine:
\begin{center}
\begin{tabular}{|l|c|l|}
\hline
memory & mean cos & reading \\
\hline
E23D01 enc block 2 & $+0.52$ & strongly aligned (conflict) \\
E23D01 enc block 3 & $+0.53$ & strongly aligned \\
trained-in D1only  & $+0.21$ & aligned \\
post-hoc \texttt{d01} & $-0.03$ & \textbf{near-orthogonal} \\
\hline
\end{tabular}
\end{center}

\end{frame}

\begin{frame}{Selection: gradient signal separates (cont.)}
On the healthy post-hoc memory, at $t=25$:
\begin{center}
\begin{tabular}{|c|c|c|c|}
\hline
$\lambda$ & mean $\|\gf\|$ sel. & mean $\|\gr\|$ sel. & overlap with AF \\
\hline
0.0 & 5.06e--3 & 1.88e--2 & 0.36 \\
1.0 & 3.79e--3 & \textbf{1.52e--3} & \textbf{0.00} \\
\hline
\end{tabular}
\end{center}

$\lambda{=}1$ gives up 25\% of forget signal to cut retain exposure
$12\times$, and selects \emph{entirely different} slots than access frequency.
\end{frame}

% ---------------------------------------------------------------------------
\begin{frame}{Ablations that came out negative (worth reporting)}
\begin{center}
\begin{tabular}{|l|c|c|c|}
\hline
post-hoc \texttt{d01} init & SH@10 & removal & UR \\
\hline
\texttt{normal} $\times 1.0$ (default) & 0.00013 & 99.0\% & \textbf{0.995} \\
\texttt{zeros}                          & 0.00009 & 99.3\% & 0.977 \\
\texttt{normal} $\times 0.01$           & 0.00004 & 99.6\% & 0.976 \\
\texttt{add} + \texttt{zeros}           & 0.00018 & 98.6\% & 0.920 \\
\hline
\end{tabular}
\end{center}

\begin{itemize}
  \item ``Smarter'' (identity-like) init is \textbf{worse}: full random init
        rebuilds best. Init \emph{scale} matters, not routing gradients.
  \item Initialising values from the FFN would \emph{re-import the spam} ---
        the removal comes precisely from discarding those weights.
\end{itemize}
\end{frame}

\begin{frame}{Ablations that came out negative (cont.)}
\begin{itemize}
  \item \textbf{Adam beat SGD} in the \texttt{pkm\_only} sweep --- but that
        updated \emph{all} slots, which is not the regime SMF's argument covers.
        Re-test under selection.
  \item \texttt{add}+\texttt{zeros} \emph{also} removes spam (98.6\%): any
        retain-only adaptation suppresses it. \alert{So we cannot easily build
        ``healthy memory $+$ spam still present''.}
\end{itemize}
\end{frame}

% ---------------------------------------------------------------------------
\begin{frame}{Encoder placements: suppression does \emph{not} stack}
Trained-in, poisoned, decoder untouched. spam$\times$ relative to non-PKM:
\begin{center}
\small
\begin{tabular}{|l|c|c|c||l|c|c|c|}
\hline
layers & N@10 & util$\times$ & spam$\times$ & layers & N@10 & util$\times$ & spam$\times$ \\
\hline
E0   & 0.03511 & 0.970 & 0.06 & E01   & 0.03511 & 0.970 & \textbf{0.04} \\
E1   & 0.03507 & 0.969 & \textbf{0.03} & E012  & 0.03315 & 0.916 & 0.24 \\
E2   & 0.03596 & 0.994 & 0.61 & E12   & 0.03717 & 1.027 & \alert{1.03} \\
E3   & 0.03508 & 0.969 & \textbf{0.03} & E23   & 0.03435 & 0.949 & \alert{2.29} \\
     &         &       &      & E0123 & 0.03739 & \textbf{1.033} & 0.26 \\
\hline
\end{tabular}
\end{center}

\begin{itemize}
  \item \textbf{No combination beats the best single layer} (E1/E3 at 0.03$\times$;
        best combo E01 at 0.04$\times$). Extra blocks buy nothing.
\end{itemize}
\end{frame}


\begin{frame}{Encoder placements: suppression does \emph{not} stack (cont.)}
\begin{itemize}
  \item \alert{Combining protective layers can destroy the protection.} E12 gives
        \emph{zero} suppression despite containing E1 (30$\times$ alone); E23
        \emph{amplifies} the attack (2.29$\times$) despite containing E3.
    \item \pierre{Pattern: E2 disrupts, E0 stabilises --- every combo with E2 but
        without E0 loses suppression; every combo with E0 keeps it. Mechanism
        unexplained.}
  \item Utility does not track layer count: E0123 (4 blocks) best (1.033),
        E012 (3 blocks) worst (0.916).
\end{itemize}
\end{frame}

% ---------------------------------------------------------------------------
\begin{frame}{PKM vs.\ fresh FFN, out of saturation: the regimes disagree}
The tie was measured where both arms hit the clean floor. Pushed out of
saturation they separate --- in \emph{opposite} directions.

\medskip
\textbf{(A) More destruction --- \texttt{dall}, 4 decoder layers:}
\begin{center}
\begin{tabular}{|l|c|c|c|}
\hline
arm & SH@10 & removal & UR \\
\hline
PKM \texttt{dall} @2000   & 0.00000 & 100\%  & 0.891 \\
PKM \texttt{dall} @10000  & 0.00000 & 100\%  & 0.686 \\
\textbf{FFN \texttt{dall} @2000}  & 0.00009 & 99.3\% & \textbf{0.964} \\
FFN \texttt{dall} @10000  & 0.00004 & 99.6\% & 0.877 \\
\hline
\end{tabular}
\end{center}
\alert{FFN wins by 7 points} and degrades less with extra steps. Where capacity
should have helped most, it \emph{hurt}.

\end{frame}

\begin{frame}{PKM vs.\ fresh FFN, out of saturation (cont.)}
\textbf{(B) Smaller budget --- \texttt{d01}, same destruction:}
\begin{center}
\begin{tabular}{|l|c|c|c|c|}
\hline
steps & PKM SH@10 & PKM UR & FFN SH@10 & FFN UR \\
\hline
200        & \textbf{0.00555} & 0.977 & 0.00863 & 0.966 \\
600        & \textbf{0.00148} & 0.989 & 0.00523 & 0.979 \\
converged  & 0.00013 & 0.995 & 0.00022 & 0.995 \\
\hline
\end{tabular}
\end{center}
\alert{PKM removes $\sim$1.75$\times$ more at every budget} and converges in
$<$half the steps (1560 vs.\ 3320).
\end{frame}

% ---------------------------------------------------------------------------
\begin{frame}{Sequential unlearning: no drift, for either method}
Same \emph{total} budget in every arm (chunks $\times$ steps $=$ 3000), same
forget set --- only the number of sequential requests varies:
\begin{center}
\begin{tabular}{|l|c|c|c|c|}
\hline
mode & chunks & steps/chunk & SH@10 & UR \\
\hline
PKM & 1  & 3000 & 0.00031 & 0.981 \\
PKM & 5  & 600  & 0.00031 & 0.992 \\
PKM & 10 & 300  & 0.00022 & 0.992 \\
FFN & 1  & 3000 & 0.00031 & \textbf{1.005} \\
FFN & 5  & 600  & 0.00022 & 0.992 \\
FFN & 10 & 300  & 0.00027 & \textbf{1.005} \\
\hline
\end{tabular}
\end{center}

\textbf{Nothing degrades.} 10 sequential requests cost the same as 1:
97.5--98.2\% removal, UR 0.981--1.005 everywhere.

\end{frame}

\begin{frame}{Sequential unlearning: no drift (cont.)}
\alert{The modular-memory hypothesis fails on its \emph{premise}:} there is no
accumulated damage for a sparse memory to protect against. A repeatedly
retrained FFN does not drift either (UR 1.005, \emph{above} clean).

\medskip
\pierre{Practical upside: deletion requests can be processed incrementally at no
cost.}
\end{frame}

% ---------------------------------------------------------------------------
\begin{frame}{Collapse is universal in trained-in PKMs}
Slot diagnostic on all five \emph{encoder-only} trained-in models:
\begin{center}
\begin{tabular}{|l|c|c|c|c|}
\hline
model & memories & slots used & Jaccard & forget-excl. \\
\hline
E01   & 2 & 128 / 128         & 1.000 & 0 \\
E012  & 3 & 128 / 95 / 32     & 1.000 & 0 \\
E12   & 2 & 125 / 32          & 1.000 & 0 \\
E23   & 2 & 32 / 32           & 1.000 & 0 \\
E0123 & 4 & 96 / 32 / 32 / 32 & 1.000 & 0 \\
\hline
\end{tabular}
\end{center}

Utilisation 0.012--0.049\% of 262{,}144 slots. \textbf{The hypothesis that the
encoder might escape collapse is refuted} --- collapse is independent of
encoder/decoder placement, layer count and depth. ($32 = knn$ exactly: those
memories return the \emph{same} 32 slots for every token.)

\end{frame}

\begin{frame}{Collapse is universal in trained-in PKMs (cont.)}
\alert{So the E2/E0 pattern cannot be a memory-content effect} --- there is no
meaningful memory in any of them. Those layers act as low-rank layers.


\end{frame}

\begin{frame}{Collapse is universal in trained-in PKMs}
\small
Slot diagnostic on all five \emph{encoder-only} trained-in models:
\begin{center}
\begin{tabular}{|l|c|c|c|c|}
\hline
model & memories & slots used & Jaccard & forget-excl. \\
\hline
E01   & 2 & 128 / 128         & 1.000 & 0 \\
E012  & 3 & 128 / 95 / 32     & 1.000 & 0 \\
E12   & 2 & 125 / 32          & 1.000 & 0 \\
E23   & 2 & 32 / 32           & 1.000 & 0 \\
E0123 & 4 & 96 / 32 / 32 / 32 & 1.000 & 0 \\
\hline
\end{tabular}
\end{center}

\medskip
\pierre{Only the post-hoc memories are healthy (59\% coverage, Jaccard 0.38,
$\sim$3000 forget-exclusive slots) $\Rightarrow$ it is not \emph{where} the memory
sits but \emph{how} it was trained: alone, against a frozen backbone.}
\end{frame}

% ---------------------------------------------------------------------------
\begin{frame}{Scorecard: PKM vs.\ a fresh FFN}
\begin{center}
\begin{tabular}{|l|l|}
\hline
regime & winner \\
\hline
Converged, single edit       & tie (saturated) \\
Heavy damage (\texttt{dall}) & \textbf{FFN} (UR 0.964 vs.\ 0.891) \\
Small budget                 & \textbf{PKM} (1.75$\times$ removal, $\sim$2$\times$ faster) \\
Sequential edits             & tie (no drift either way) \\
\hline
\end{tabular}
\end{center}

\medskip
\textbf{PKM's only measured advantage is convergence speed at constrained
budgets.} The post-hoc contribution is therefore the \emph{localised retrain}
itself --- re-initialise the first two decoder FFNs, retrain on retain data,
backbone frozen --- with PKM a speed optimisation on top, not the mechanism.
\end{frame}

% ---------------------------------------------------------------------------
\begin{frame}{Fixing collapse: what the literature actually says}
\textbf{Neither paper offers a recipe --- and the reason matters.}

\medskip
\textbf{Sparse Memory Finetuning} (\url{https://arxiv.org/abs/2510.15103}) never
faced collapse: it \emph{inherits} a memory pretrained at scale.
\begin{center}
\begin{tabular}{|l|l|}
\hline
their setup & 1.3B base, 1M slots, $k{=}32$, 4 heads, $d_v{=}1024$ \\
placement & ONE layer, middle of the model (12 of 22) \\
optimizer & \textbf{SGD, lr 0.1--2} (AdamW for baselines) \\
selection & TF-IDF over 1000 background batches, $t \in \{25..1000\}$ \\
\hline
\end{tabular}
\end{center}
No initialisation scheme, no warm-up, no query normalisation, \alert{no slot
utilisation diagnostics at all}.

\end{frame}

\begin{frame}{Fixing collapse: what the literature says (cont.)}
\textbf{Memory Layers at Scale} (2412.09764), the source of that layer, mentions
only \textbf{qk-normalisation} ``when needed'' for instability, plus silu gating
and a \emph{shared} pool across 3 layers. No load-balancing loss.

\medskip
\pierre{So the collapse measurement --- and any fix --- is ours. Our lr is
$10^{-3}$; SMF uses 0.1--2, i.e.\ $100\text{--}1000\times$ larger.}
\end{frame}

% ---------------------------------------------------------------------------
\begin{frame}{Two candidate fixes, and why collapse self-reinforces}
\textbf{The trap.} At init the values are random, so \emph{which} slots the router
picks barely changes the loss $\Rightarrow$ almost no gradient pressure on the
query net $\Rightarrow$ it collapses to a near-constant. After that only the
surviving slots ever train, which reinforces the collapse.

\end{frame}

\begin{frame}{Candidate fix 1: batchnorm on the query}
\textbf{(1) \texttt{query\_norm=batchnorm}} (\texttt{track\_running\_stats=False},
so the padding-corrupted running stats that motivated disabling it never exist):
\begin{center}
\begin{tabular}{|l|c|c|}
\hline
query\_norm & token-shared $\|\mu\|^2$ & distinct top-$k$ patterns \\
\hline
none      & 144.45 & 59 / 2048 \\
layernorm & 12.56  & 51 / 2048 \quad(\emph{worse}) \\
\textbf{batchnorm} & \textbf{0.000} & \textbf{118 / 2048} \\
\hline
\end{tabular}
\end{center}
Zeroes the token-shared component exactly and doubles diversity ---
\alert{but still concentrated, so only a PARTIAL fix.}

\end{frame}

\begin{frame}{Candidate fix 2: warm-up noise}
\textbf{(2) \texttt{warmup\_noise}} --- noisy top-$k$ gating: perturb routing
scores for the first $N$ steps so selection varies and gradient spreads, at
\emph{identical} cost. On a toy that reproduces collapse:
\begin{center}
\begin{tabular}{|l|c|}
\hline
config & slots touched / 4096 \\
\hline
baseline  & \alert{33} \quad(0.8\%, collapse reproduced) \\
noise 0.5 & 1776 \\
\textbf{noise 2.0} & \textbf{4082} \quad(99.7\%) \\
\hline
\end{tabular}
\end{center}
\end{frame}

% ---------------------------------------------------------------------------
\begin{frame}{Collapse fix: status --- \alert{verdict NOT in yet}}
\textbf{What is measured:} only the two synthetic checks on the previous slide
(query geometry; slot spread \emph{during} the warm-up window on a toy module).

\medskip
\alert{What is NOT measured:} whether either knob raises slot utilisation in an
actually-trained TIGER model above the 0.012--0.049\% seen in all 8 layouts ---
and critically, \textbf{whether utilisation survives after the noise switches
off.} The router may re-collapse once selection becomes deterministic again. The
toy cannot answer this.

\end{frame}

\begin{frame}{Collapse fix: status (cont.)}
In flight: 4 capped runs on E1 (\texttt{max\_steps=4000}), then the slot
diagnostic on each.
\begin{center}
\begin{tabular}{|l|l|c|}
\hline
arm & config & status \\
\hline
base  & ---                                 & \strun \\
bn    & \texttt{query\_norm=batchnorm}       & \strun \\
wu    & \texttt{warmup\_noise=2.0}, 2000 st. & \sttodo \\
bnwu  & both                                 & \sttodo \\
\hline
\end{tabular}
\end{center}

\end{frame}

\begin{frame}{Collapse fix: status (cont.)}
\pierre{The first warm-up attempt raised \texttt{knn} instead of adding noise.
Lookup cost scales as $knn^2$, so $knn{=}256$ built a $(\text{rows},65536)$
tensor, overflowed int32 indexing in \texttt{topk} and aborted with ``CUDA error:
illegal memory access''. \texttt{warmup\_knn} is now hard-capped at 64.}
\end{frame}

% ===========================================================================
\section{Plan}

\begin{frame}{What is done}
\begin{center}
\begin{tabular}{|l|c|}
\hline
item & status \\
\hline
PKM in training / post-hoc insertion            & \stdone \\
Memory-only update (\texttt{pkm\_only})          & \stdone \\
Full retain stream (\texttt{retain\_source})     & \stdone \\
Slot access counters + diagnostics               & \stdone \\
Top-$t$ selection (AF / AF-IHF / grad / grad+dot)& \stdone \\
Optimizer + init knobs                           & \stdone \\
\hline
\end{tabular}
\end{center}
\end{frame}

\begin{frame}{What is done (cont.)}
\begin{center}
\begin{tabular}{|l|c|}
\hline
item & status \\
\hline
Anti-collapse knobs (\texttt{query\_norm}, \texttt{warmup\_noise}) & \stdone \\
FFN-reinit control (\texttt{ffn\_reinit\_layers})  & \stdone \\
Sequential-unlearning comparison                  & \stdone \\
Placement screen (post-hoc, 8 layers)            & \stdone \\
Placement screen (trained-in, 8 layers)          & \stdone \\
Architecture study + 168-run unlearning sweeps   & \stdone \\
Clean-baseline audit (seed noise)                & \stdone \\
\hline
\end{tabular}
\end{center}
\end{frame}

% ---------------------------------------------------------------------------
\begin{frame}{What needs to be done}
\begin{center}
\begin{tabular}{|l|l|c|}
\hline
\# & item & status \\
\hline
1 & \textbf{FFN-reinit control} instead of PKM replace   & \stdone \\
2 & Other unlearning algorithms under \texttt{pkm\_only} & \sttodo \\
3 & Fix memory collapse (routing, not capacity)          & \stdone \\
4 & Top-$t$ sweep on healthy (trained-in) memory         & \stdone \\
5 & AF-IHF vs.\ grad vs.\ access baseline ablation       & \stdone \\
6 & Adam vs.\ SGD \emph{under} selection                 & \sttodo \\
7 & Why does trained-in plateau at SH $\approx$ 0.004?   & \sttodo \\
8 & Multi-seed (clean baselines \emph{and} runs)         & \strun \\
9 & Other strategies $+$ other (Amazon) datasets         & \sttodo \\
\hline
\end{tabular}
\end{center}

\end{frame}

\begin{frame}{What needs to be done (cont.)}
\alert{(1) is the sanity check the post-hoc story rests on:} re-initialise the
same two decoder FFNs as \emph{ordinary FFNs} and fine-tune them identically.
If that matches PKM, the mechanism is ``reinit-and-retrain-this-layer'' and the
PKM is an implementation detail.

\medskip
\textbf{(3) is blocking (4)--(6).} Candidate causes:
\texttt{query\_batchnorm=false} (padding-unsafe), \texttt{query\_proj} collapsing
to a constant, keys never leaving their init.
\pierre{Cheapest lead: post-hoc PKMs do not collapse --- isolate the memory
during training.}
\end{frame}

% ---------------------------------------------------------------------------
\begin{frame}{Item 1: the FFN-reinit control (does PKM actually help?)}
The post-hoc recipe is: \emph{discard} a trained FFN, put a fresh high-capacity
module there, fine-tune it on retain data only.

\medskip
\textbf{Control:} do exactly the same, but re-initialise the layer as an
\emph{ordinary FFN} --- same two decoder blocks, same discovery code, same
optimiser, same steps, same full retain split.
\begin{center}
\begin{tabular}{|l|c|c|c|}
\hline
arm & SH@10 & removal & UR \\
\hline
poison baseline               & 0.01274 & ---    & 0.962 \\
PKM \texttt{d01} $+$ rebuild  & 0.00013 & 99.0\% & 0.995 \\
\textbf{fresh FFN \texttt{d01} (control)} & \textbf{0.00022} & \textbf{98.2\%} & \textbf{0.995} \\
\hline
\end{tabular}
\end{center}

\end{frame}

\begin{frame}{Item 1: the FFN-reinit control (cont.)}
\alert{Why this decides the framing:} the removal already appears to come from
the \emph{retain-only rebuild}, not from the memory --- ablation alone removes
only 12\%, and \texttt{add} mode (FFN left intact) still removes 98.6\%.
\alert{Result: tied.} Utility identical to 4 decimals (0.03745 vs.\ 0.03746),
removal 99.0\% vs.\ 98.2\%. A plain FFN at \textbf{1/25th} the parameter budget
(2.6M vs.\ 67M) does it just as well.

\end{frame}

\begin{frame}{Item 1: the FFN-reinit control (cont.)}
\textbf{So the headline post-hoc result is a \emph{localised retrain}, not a
memory method} --- and it needs no PKM machinery at all.

\medskip
\pierre{But the test is SATURATED: both arms sit at the clean floor. A tie under
saturation shows this task cannot separate them, not that they are equivalent.}
\pierre{Implemented as \texttt{model.ffn\_reinit\_layers} $+$
\texttt{unlearning.update\_scope=ffn\_only}; the reinit runs AFTER the ckpt load
(FFN weights are in the ckpt, PKM's are not).}
\end{frame}

% ---------------------------------------------------------------------------
\begin{frame}{Item 7: what is the ``SH $\approx$ 0.004 plateau''?}
On the trained-in E23D01 model \emph{every} gradient-based route stops at the
same residual, regardless of parameter scope or budget:
\begin{center}
\begin{tabular}{|l|c|c|}
\hline
route & SH@10 & UR \\
\hline
\texttt{unified}, full model   & 0.00411 & 1.013 \\
\texttt{pkm\_only}, 40 epochs  & 0.00452 & 1.017 \\
\texttt{pkm\_only}, 100 epochs & 0.00385 & 1.004 \\
\texttt{pkm\_only}, 200 epochs & 0.00496 & 1.011 \\
\hline
\end{tabular}
\end{center}

\textbf{Not a measurement floor:} this model's \emph{clean} SH@10 is 0.00004 ---
the plateau sits $\sim$85$\times$ above it, so real spam remains.
\textbf{Not irreducible:} post-hoc reaches 0.00013 on the same attack.

\end{frame}

\begin{frame}{Item 7: the plateau (cont.)}
Three candidate explanations, each with a different implication:
\begin{enumerate}
  \item \textbf{Optimisation} --- more/better steps would get there.
        \pierre{Weakened: 40$\to$100$\to$200 epochs is non-monotone.}
  \item \textbf{Redundancy} --- spam spread over parameters the objective cannot
        move simultaneously.
  \item \textbf{It is the utility-preserving frontier}, not a ceiling: pushing SH
        lower costs utility. \pierre{\texttt{seif} reaches SH$\approx$0 at
        UR 0.67--0.74, consistent with this.}
\end{enumerate}

\medskip
\alert{If (3), ``ceiling'' is the wrong word} --- the question becomes why the
post-hoc rebuild escapes the frontier entirely.
\end{frame}

% ---------------------------------------------------------------------------
\begin{frame}{Open design question}
The selection experiment needs a model with
\textbf{(a)} healthy memory \emph{and} \textbf{(b)} spam still present.

\medskip
But every route to (a) currently destroys (b):
\begin{itemize}
  \item \texttt{replace} + rebuild $\Rightarrow$ healthy, spam gone (99\%)
  \item \texttt{add} + warmup $\Rightarrow$ healthy, spam gone (98.6\%)
  \item trained-in $\Rightarrow$ spam present, memory collapsed
\end{itemize}

\end{frame}

\begin{frame}{Open design question (cont.)}
\textbf{Options:}
\begin{enumerate}
  \item Fix collapse in trained-in models (then (a) and (b) coexist naturally)
  \item Two-stage: warm \emph{briefly} (routing forms, spam survives)
        $\rightarrow$ select $\rightarrow$ train selected slots only
  \item Reframe: selection as a way to make the \emph{rebuild} cheaper/safer
        rather than to target spam
\end{enumerate}

\medskip
\todo{decide which framing the paper takes}
\end{frame}

% ---------------------------------------------------------------------------
\begin{frame}{Threats to validity}
\begin{itemize}
  \item \textbf{Single seed everywhere.} Three identical-config clean runs span
        \textbf{7.5\%} (sd 3.6\%) --- larger than most architecture effects.
        Utility claims below $\sim$5\% are not currently supportable.
  \item \textbf{One attack, one dataset, one target strategy}
        (beauty / bandwagon / \texttt{mid}).
  \item \textbf{Denominators.} UR must be taken against each model's \emph{own}
        clean baseline; mixing them conflates architecture cost with unlearning
        cost.
  \item \textbf{One cell has no attack} (ALLv2/popular: poison 0.00018 vs.\
        clean 0.00013) --- must be excluded from removal claims.
\end{itemize}
\end{frame}

% ---------------------------------------------------------------------------
\begin{frame}{Resources}
\begin{itemize}
  \item ERASE: \url{https://arxiv.org/abs/2603.08341}
  \item FFNs are key--value memories: \url{https://aclanthology.org/2021.emnlp-main.446/}
  \item Product-key memory (Lample et al.\ 2019): \url{https://arxiv.org/abs/1907.05242}
  \item Memory layers at scale: \url{https://github.com/facebookresearch/memory}
  \item GRID: \url{https://github.com/snap-research/GRID}
\end{itemize}

\medskip
\textbf{Internal.} \texttt{WORKFLOW.md} sections F--J (full experiment record,
job ids); \texttt{tables/pkm\_update.md} (extended notes).

\medskip
\textbf{Deadline.} ICLR 2027 --- Sep 25, 2026 AOE.
\end{frame}




\begin{frame}{PKM as adapter}
\begin{itemize}
    \item Remember: 4 encoder, 4 decoder layers in TIGER
    \item PKM as adapter after training is better in the decoder layers than encoder layers empirically from previous experiments
    \item When replacing FFNs in all decoder layers, FFN replacement performs better then PKM replacement but PKM is faster
    \item but when replacing only FFNs of the decoders with index 0 and 1 then PKM performs better and faster


\end{itemize}
\end{frame}

% ---------------------------------------------------------------------------
\begin{frame}{PKM vs.\ fresh FFN: PKM buys speed, not final quality}

\textbf{The comparison depends on the available retraining budget.}

\medskip
\textbf{(A) Small budget: PKM converges substantially faster}

\begin{center}
\begin{tabular}{|c|c|c|c|c|}
\hline
steps & PKM SH@10 & PKM UR & FFN SH@10 & FFN UR \\
\hline
200 & \textbf{0.00555} & \textbf{0.977} & 0.00863 & 0.966 \\
600 & \textbf{0.00148} & \textbf{0.989} & 0.00523 & 0.979 \\
\hline
converged &
0.00013 & 0.995 &
0.00022 & 0.995 \\
\hline
\end{tabular}
\end{center}

PKM removes $\sim$1.75$\times$ more spam at \emph{both} intermediate
budgets and reaches convergence in 1560 vs.\ 3320 steps.


\end{frame}


\begin{frame}{PKM vs.\ fresh FFN: PKM buys speed, not final quality}

\textbf{(B) Heavy damage: FFN is better}

\begin{center}
\begin{tabular}{|l|c|c|c|}
\hline
arm & SH@10 & removal & UR \\
\hline
PKM \texttt{dall} @2000 & 0.00000 & 100\% & 0.891 \\
\textbf{FFN \texttt{dall} @2000} & \textbf{0.00009} & \textbf{99.3\%} & \textbf{0.964} \\
PKM \texttt{dall} @10000 & 0.00000 & 100\% & 0.686 \\
FFN \texttt{dall} @10000 & 0.00004 & 99.6\% & 0.877 \\
\hline
\end{tabular}
\end{center}

\alert{PKM's advantage is convergence speed, not final quality.}

\end{frame}

% ---------------------------------------------------------------------------
\begin{frame}{PKM is a better adapter for localized repair}

\textbf{Post-hoc d01:} replace the first two decoder FFNs and retrain on
retain data only, with the backbone frozen.

\medskip

\begin{center}
\begin{tabular}{|c|c|c|c|c|}
\hline
steps & PKM SH@10 & PKM UR & FFN SH@10 & FFN UR \\
\hline
200 &
\textbf{0.00555} & \textbf{0.977} &
0.00863 & 0.966 \\
600 &
\textbf{0.00148} & \textbf{0.989} &
0.00523 & 0.979 \\
\hline
converged &
\textbf{0.00013} & 0.995 &
0.00022 & 0.995 \\
\hline
\end{tabular}
\end{center}


\end{frame}


\begin{frame}{PKM is a better adapter for localized repair (cont.)}
\textbf{PKM reaches the same final solution with much less optimization:}

\begin{itemize}
  \item At 200 steps: \textbf{1.55$\times$ more removal}
        ($56.5\%$ vs.\ $32.3\%$).
  \item At 600 steps: \textbf{1.50$\times$ more removal}
        ($88.4\%$ vs.\ $58.9\%$).
  \item Convergence: \textbf{1560 vs.\ 3320 steps} --- less than half.
  \item Final utility is essentially identical: UR $=0.995$.
\end{itemize}

\medskip

\alert{\textbf{For the localized d01 adapter, PKM is a substantially faster
way to relearn the destroyed function.}}

\medskip

\textbf{Interpretation:} sparse top-$k$ retrieval updates only the relevant
memory rows, allowing the adapter to fit the retain distribution rapidly.

\end{frame}


\begin{frame}{Collapse for trained-in PKM}

\begin{itemize}
    \item Problem: Memory collapse when training with PKM from the start
    \item Batch norm can increase memory utilization from experiments (0.01\% used to 1\% used)
\end{itemize}

\end{frame}





% ===========================================================================
\section{Trained-in PKM: fixing the collapse (Oct 2026)}

\begin{frame}{New setup for the trained-in experiments}
\begin{itemize}
  \item TIGER with a PKM \emph{replacing} the FFN of \textbf{encoder block 1}
        (E1), trained from scratch on the poisoned data (beauty, bandwagon
        \texttt{mid}, 1\%), 4000 steps, deterministic, RQ-VAE semantic IDs.
  \item Collapse is judged with the \textbf{eval-mode slot diagnostic}: slots read
        by forget / retain data, Jaccard, forget-exclusive slots.
  \item New training-time logging: live slots per window (a trend only: dropout
        0.15 makes it overstate eval routing $\sim$10$\times$).
\end{itemize}
\end{frame}

\begin{frame}{New setup: own reference points}
Same setup, test set:
\begin{center}
\begin{tabular}{|l|c|c|}
\hline
model & N@10 & SH@10 \\
\hline
poisoned, no fix (p0)       & 0.0307 & 0.0022 \\
clean, no fix               & 0.0321 & 0.0001 \\
clean, with balance fix     & 0.0313 & 0 \\
\hline
\end{tabular}
\end{center}
\end{frame}

% ---------------------------------------------------------------------------
\begin{frame}{Attempt 1: explore / exploit dense updates}
With probability $p$ per step every PKM layer reads (and updates)
\emph{all} $512^2$ slots, otherwise the usual top-$k$. Exact and cheap:
$\text{softmax}(s_1{+}s_2) = \text{softmax}(s_1)\otimes\text{softmax}(s_2)$.
Schedules: constant, linear, cosine, adaptive (lower $p$ as utilisation rises).

\medskip
\begin{center}
\small
\begin{tabular}{|l|c|c|c|}
\hline
E1 arm & live slots @1000 & @2000 & @4000 \\
\hline
$p=0$ (baseline)         & 16.9k & 1.3k & \alert{32} \\
dense $p=0.05$           & 15.8k & 1.7k & 818 \\
dense $p=0.2$            & 15.7k & 1.5k & \alert{32} \\
dense $p=0.5$            & 60.8k & 32   & \alert{32} \\
straight-through $p=0.2$ & 15.5k & 2.1k & 1171 \\
adaptive / linear        & 45--62k & 32 & \alert{32} \\
\hline
\end{tabular}
\end{center}
\end{frame}

\begin{frame}{Attempt 1: explore / exploit --- verdict: negative}
\begin{itemize}
  \item \textbf{First measurement of collapse dynamics:} the memory starts fully
        used and the router narrows over $\sim$2000 steps, in encoder \emph{and}
        decoder.
  \item Dense updates do \textbf{not} stop it; high $p$ collapses
        \emph{faster} (32 slots by step 1500). The dense gradient rewards sub-keys
        that are good \emph{on average}, a token-independent direction.
  \item Eval diagnostic: every arm reads 32--160 slots, 0 forget-exclusive.
\end{itemize}

\medskip
\pierre{Collapse is a \textbf{routing} problem: giving every value a gradient does
not change the router's incentive.}
\end{frame}

% ---------------------------------------------------------------------------
\begin{frame}{Attempt 2: a routing balance loss}
Added to the training loss, per head and sub-key half ($n{=}512$ sub-keys):

\medskip
\textbf{MI form:} $\;\log n - H(\bar p) + \tfrac1B\sum_b H(p_b)$
\quad ($p_b$ = token $b$'s sub-key softmax)
\begin{itemize}
  \item the plain marginal term is fooled by identical flat routing (= collapse)
  \item measured: \alert{inert} at small coef (flat scores, zero gradient),
        \alert{sharpens tokens} at large coef. Not scale-invariant.
\end{itemize}
\end{frame}

\begin{frame}{Attempt 2: a routing balance loss (switch form)}
\textbf{Switch form} (load balancing on the \emph{hard} assignment):
\;$n\sum_i f_i P_i$
\begin{itemize}
  \item $f_i$ = share of top-$k$ assignments to sub-key $i$ (counted)
  \item $P_i$ = mean softmax on \textbf{per-token standardised} scores, so it
        cannot be satisfied by shrinking all scores
\end{itemize}
\end{frame}

\begin{frame}{Attempt 2: the switch loss fixes the collapse}
Eval slot diagnostic, E1, retain subset (6.6k sequences):
\begin{center}
\small
\begin{tabular}{|l|c|c|c|c|}
\hline
arm & retain slots & Jaccard & forget-excl. & test N@10 \\
\hline
p0 (no fix)              & \alert{32} & 1.00 & 0 & 0.0307 \\
batchnorm only           & 2.5k  & 0.71 & 107  & 0.0303 \\
MI, coef 0.1             & 2.4k  & 0.84 & 65   & 0.0328 \\
MI + bn, coef 0.1        & 5.6k  & 0.81 & 121  & 0.0299 \\
switch, coef 0.1         & 63k   & 0.68 & 2.6k & 0.0322 \\
\textbf{switch + bn, coef 0.1} & \textbf{92k} & \textbf{0.43} & 16.7k & 0.0305 \\
\hline
\end{tabular}
\end{center}

\begin{itemize}
  \item 2900$\times$ the slots of p0 at \textbf{unchanged utility}.
  \item \alert{Full retain set (22k seqs):} 165k slots (63\%), Jaccard 0.39, but
        only \textbf{98} forget-exclusive (the 16.7k was a sampling artefact).
\end{itemize}
\end{frame}

\begin{frame}{The fix replicates}
Eval diagnostic (forget-exclusive on the full retain set where given):
\begin{center}
\small
\begin{tabular}{|l|c|c|c|}
\hline
run & retain slots & Jaccard & test N@10 (p0) \\
\hline
seed 2              & 165k (63\%) & 0.39 & 0.0305 (0.0307) \\
seed 3              & 147k (56\%) & 0.40 & 0.0320 (0.0318) \\
seed 4              & 123k (47\%) & 0.39 & 0.0328 (0.0326) \\
coef 0.3            & 193k (74\%) & 0.41 & 0.0301 \\
decoder D01         & 21k / 24k\,$^*$ & 0.58 & 0.0299 \\
sports 1\%          & 56k\,$^*$   & 0.51 & 0.0188 (0.0209) \\
\hline
\end{tabular}
\end{center}
{\footnotesize $^*$ retain subset.}

\begin{itemize}
  \item Forget-exclusive slots stay $\approx$0 (49--98) at every seed:
        selection must use frequency \emph{contrast}, not exclusivity.
  \item \alert{Sports}: $-10\%$ N@10 (one seed).
\end{itemize}
\end{frame}

% ===========================================================================
\section{Unlearning with a healthy memory}

\begin{frame}{A healthy memory is required for memory-only unlearning}
\texttt{unified}, memory values only (backbone frozen), recipe of the best
recorded \texttt{pkm\_only} run, 100 epochs:
\begin{center}
\small
\begin{tabular}{|l|c|c|}
\hline
arm & N@10 & SH@10 \\
\hline
balanced, before unlearning        & 0.0305 & 0.0052 \\
balanced, all values               & 0.0286 & 0.00004 \\
balanced, AF-IHF top $t{=}100$     & 0.0277 & \textbf{0} \\
balanced, grad+dot top $t{=}1000$  & 0.0283 & 0 \\
\hline
collapsed p0, before               & 0.0307 & 0.0022 \\
collapsed p0, memory only          & 0.0294 & \alert{0.0090} \\
\hline
\end{tabular}
\end{center}
\end{frame}

\begin{frame}{A healthy memory is required: takeaways}
\begin{itemize}
  \item Editing \textbf{100 of 262k value rows} removes the spam completely.
  \item On the collapsed memory the same update \alert{increases spam 4$\times$}
        (2 of 3 seeds).
  \item At 100 epochs the test is \emph{saturated} (SH 0 everywhere).
\end{itemize}
\end{frame}

\begin{frame}{Budget sweep: AF-IHF beats grad+dot}
N@10 / SH@10 after unlearning (before: 0.0305 / 0.0052):
\begin{center}
\small
\begin{tabular}{|l|c|c|c|}
\hline
arm & 4 epochs & 10 epochs & 40 epochs \\
\hline
all values         & .0290 / .0001 & .0280 / 0 & .0283 / 0 \\
AF-IHF $t{=}100$   & .0303 / .0095 & \textbf{.0291 / .0003} & .0277 / 0 \\
grad+dot $t{=}100$ & .0305 / .0104 & .0303 / .0049 & \textbf{.0291 / .0002} \\
AF-IHF $t{=}1000$  & \textbf{.0294 / .0005} & .0280 / 0 & .0277 / 0 \\
grad+dot $t{=}1000$& .0304 / .0061 & .0296 / .0014 & .0282 / .00004 \\
\hline
\end{tabular}
\end{center}

\begin{itemize}
  \item AF-IHF removes much faster than grad+dot at equal budget.
  \item Best selection arms: 91--96\% removal at $-3.6$ to $-4.6\%$ N@10;
        editing all values costs $-8\%$ for full removal.
\end{itemize}
\end{frame}

\begin{frame}{But the full model still wins on spam}
Full-model \texttt{unified} with \emph{its own} recorded recipe (4 epochs),
UR against the own \textbf{clean} model (0.0313):
\begin{center}
\small
\begin{tabular}{|l|c|c|c|}
\hline
method & seed 2 & seed 3 & seed 4 \\
\hline
full model           & \textbf{.0313} (UR 1.00) & \textbf{.0324} & \textbf{.0342} \\
best memory-only     & .0294 (UR 0.94) & .0311 & .0320 \\
\hline
\end{tabular}
\end{center}
All arms reach SH@10 $\approx$ 0.

\begin{itemize}
  \item Gap of \textbf{4--6\% N@10 at every seed}; the full model matches the
        retrained model exactly.
  \item \textbf{50 sequential requests} instead of 10: every method loses
        3--4\%, memory-only does \emph{not} accumulate less damage.
  \item Memory-only's remaining argument: \emph{locality} (100--1000 value rows,
        13k--128k parameters change).
\end{itemize}
\end{frame}

\begin{frame}{Sensitive items: setting too small to decide}
Beauty, hair-loss products, forget ratio $10^{-4}$ (14 users), per-user deletion.
SHF / SHR = sensitive hits for requesting / other users.
\begin{center}
\small
\begin{tabular}{|l|c|c|c|}
\hline
balanced model & N@10 & SHF@10 & SHR@10 \\
\hline
retrained (target)   & 0.0316 & 0.357 & 0.0378 \\
before unlearning    & 0.0301 & 0.214 & 0.0378 \\
full model           & 0.0294 & 0.071 & \alert{0.0239} \\
values only          & 0.0303 & 0.143 & 0.0367 \\
AF-IHF $t{=}100$     & 0.0306 & 0.214 & 0.0379 \\
\hline
\end{tabular}
\end{center}

\begin{itemize}
  \item The retrained model \emph{is} the original here: deleting 23
        interactions changes nothing measurable.
  \item The full model \alert{over-deletes for everyone} (SHR $-37\%$).
        Memory-only stays close, but only by changing little.
  \item \todo{repeat with the $10^{-3}$ forget ratio}
\end{itemize}
\end{frame}

% ===========================================================================
\section{MUNKEY-inspired extensions (in progress)}

\begin{frame}{MUNKEY: models designed to forget}
\textbf{Laguna et al.\ 2026}, \emph{Machine Unlearning via Key Deletion}
(\url{https://arxiv.org/abs/2603.15033}). Domain: vision.

\medskip
\begin{itemize}
  \item \emph{Unlearning by design}: a memory-augmented ViT separates
        instance-specific memorisation from the model weights.
  \item Every training instance $x_i$ gets a key $k_i = g_\phi(x_i)$ from a
        \textbf{frozen} encoder and a \textbf{learnable} exemplar token $v_i$.
  \item Forgetting $=$ deleting the pair $(k_f, v_f)$: zero-shot, no weight
        update, no access to the data. Instance-level unlearning.
\end{itemize}

\medskip
\textit{``By repurposing memory-augmented structures for the `inverse' task of
forgetting, we enable a new class of models that are not only high-performing
but are natively compliant with the dynamic requirements of data privacy and
integrity.''}
\end{frame}

\begin{frame}{MUNKEY: architecture}
\begin{itemize}
  \item $g_\phi$: any suitable visual encoder, \textbf{pretrained} (ViT-Base
        on ImageNet by default).
  \item Adapter $h_\psi$ (a single linear layer) projects the exemplar token
        into the ViT latent space: $v_i' = h_\psi(v_i)$, used as an auxiliary
        input token next to the image patch embeddings $z_i$.
  \item $v_i$ is \emph{not} derived from $k_i$; the key is only used for
        retrieval.
\end{itemize}

\medskip
\pierre{Pretrained $g_\phi$: what if the forgotten images are in its
pretraining data?}

\medskip
\pierre{Memory overhead is linear in dataset size (one key and one token per
instance). Manageable here, but scalability (retrieval latency, training
overhead, $K$ forward passes at inference) is unverified.}
\end{frame}

\begin{frame}{MUNKEY: training with pathway dropout}
Stochastic mask over the two pathways (image tokens, exemplar token):
\begin{gather*}
  m = [\gamma_{img}, \gamma_{tok}] \in \{0,1\}^2 \\
  \mathbb{P}(m{=}[0,1]) = p_i,\quad
  \mathbb{P}(m{=}[1,0]) = p_t,\quad
  \mathbb{P}(m{=}[0,0]) = 0
\end{gather*}
and $\mathbb{P}(m{=}[1,1]) = 1 - p_i - p_t$. A dropped pathway is replaced by a
globally learned null token:
\[
  D(u, \emptyset, \gamma) = \gamma u + (1-\gamma)\,\emptyset
\]

\medskip
\pierre{A learned null token for a dropped pathway: does it end up holding
general knowledge?}
\end{frame}

\begin{frame}{MUNKEY: retrieval-based regularisation}
At inference the query's own token is unavailable, so this state is simulated in
training: with probability $p_r = 0.2$, use the similarity-weighted average
$v_n'$ of $K' \sim U\{2,\ldots,16\}$ retrieved neighbour tokens (self excluded):
\[
  z_{in}(x_i, m, r) = \big[\,D(z_i, \emptyset_{img}, \gamma_{img}),\;
  D\big((1-r)\,v_i' + r\,v_n',\ \emptyset_{tok}, \gamma_{tok}\big)\big]
\]
$\theta$, $\psi$ and all $(v_i)_{i=1}^N$ are optimised jointly on the expected
masked objective.

\medskip
\begin{itemize}
  \item Gradients into neighbour tokens are \textbf{clipped to 0} when the
        retrieval branch is active, so the backbone learns to aggregate context
        and neighbours do not drift towards one sample's gradients.
\end{itemize}
\end{frame}

\begin{frame}{MUNKEY: inference and unlearning}
\textbf{Inference.} For a query $x_q$, $k_q = g_\phi(x_q)$ retrieves the $K$
nearest keys. The model runs \emph{once per neighbour} on $[z_q, v_j']$, and the
outputs are ensembled with weights
$w_j = \text{softmax}_j\big(\cos(k_q, k_j)/\tau\big)$.

\medskip
\textbf{Unlearning.} Delete $(k_f, v_f)$ for each forget sample.

\medskip
\pierre{How is it ensured this has unlearned? Could the global null tokens
retain knowledge?}
\end{frame}

\begin{frame}{MUNKEY: pathway sensitivity and model selection}
Guard against \emph{pathway collapse}: the model ignores the memory (poor
unlearning) or degenerates into a lookup table (poor generalisation).

\medskip
\begin{itemize}
  \item Sensitivity score $P_s$: the accuracy with the image pathway only minus
        the accuracy with the memory pathway only, divided by the accuracy with
        both (plus a stability term). The other pathway is replaced by its null
        token.
  \item Multi-stage hyperparameter selection: choose $(p_i, p_t)$ with
        $P_s(p_i, p_t) < \xi$ for a threshold $\xi \ge 0$, then the smallest
        validation gap.
\end{itemize}
\end{frame}

\begin{frame}{MUNKEY: pathway sensitivity and model selection (cont.)}
\begin{itemize}
  \item Null tokens give pathway-specific accuracy without unfamiliar inputs
        (the model saw them in training).
\end{itemize}

\medskip
\pierre{Why should a low $P_s$ imply good forgetting?}
\end{frame}

\begin{frame}{MUNKEY: experimental setup}
\begin{itemize}
  \item \textbf{Natural images:} CIFAR-10, CIFAR-100, Tiny ImageNet.
  \item \textbf{Medical images:} BloodMNIST, PathMNIST, DermaMNIST (MedMNIST v2).
  \item \textbf{Utility:} test accuracy and retain accuracy.
        \textbf{Forgetting:} forget accuracy and MIA.
  \item MIA: RMIA (\url{https://arxiv.org/abs/2312.03262}) with 64 shadow models.
  \item Random forget sets of 10\% and 2\%, stratified to keep the label
        distribution; 3 seeds; component ablations in their appendix D.
\end{itemize}

\medskip
\pierre{Why retain accuracy? How does performance compare to the baselines
\emph{before} unlearning?}
\end{frame}

\begin{frame}{MUNKEY: results}
\begin{itemize}
  \item \textit{``The key encoder does not constitute an unfair advantage as each
        method's unlearning performance is evaluated against its respective gold
        standard, i.e.\ Retrain (T) or Retrain (M).''}
  \item \pierre{Retrain (M) retrains the MUNKEY architecture, Retrain (T) the
        tiny ViT. MUNKEY uses frozen ViT-Base keys, $>$10$\times$ the parameters.
        The gain may come from the better key model: if the best model were used
        everywhere, the key choice would lose that advantage.}
\end{itemize}
\end{frame}

\begin{frame}{MUNKEY: results (cont.)}
\begin{itemize}
  \item \textbf{Memory content:} the ViT CLS token quickly forms class-specific
        clusters, while exemplar tokens evolve more granularly (\pierre{one token
        per sample?}).
  \item Token-level aggregation (softmax / rank, one forward pass) is called
        competitive when $K$ passes are too expensive.
        \pierre{Their Table 3 shows a higher average gap.}
\end{itemize}
\end{frame}

\begin{frame}{MUNKEY: questions and answers (1/3)}
\begin{itemize}
  \item \pierre{Ensembling usually adds a lot. Compared to an ensembled ViT?}
        No. With $K{=}4$ they only compare token-level softmax and rank
        aggregation, deep prompting and cross-attention.
  \item \pierre{What are the null tokens for?} An evaluation device: they let
        $P_s$ be measured. They may capture general knowledge, but are not
        designed to.
  \item \pierre{What is $h_\psi$?} A single linear layer mapping the
        independently learned $v_i$ into the ViT embedding space.
  \item \pierre{Seeds?} Three per experiment.
\end{itemize}
\end{frame}

\begin{frame}{MUNKEY: questions and answers (2/3)}
\begin{itemize}
  \item \pierre{Why would $\theta$ not contain knowledge about the samples?}
        It does; the removal is an empirical claim. Possibly it works because
        the neighbours' signal in embedding space dominates: the model treats
        the sample as part of its neighbourhood. That is neighbourhood
        unlearning done by the architecture rather than by an algorithm, in the
        well-trained space of the frozen encoder.
  \item This favours random-instance forgetting: a close, heavily weighted
        neighbour usually survives the deletion.
\end{itemize}
\end{frame}

\begin{frame}{MUNKEY: questions and answers (3/3)}
\begin{itemize}
  \item \pierre{Which MIA, too weak?} RMIA is a strong black-box attack, but
        no adaptive or white-box attacks are tested, and membership in the
        frozen key encoder is ignored.
  \item \pierre{Multiple requests?} None sequentially (deletion makes it
        trivial); no hard vs.\ easy forget sets.
\end{itemize}
\end{frame}

\begin{frame}{MUNKEY: criticism (1/3)}
\begin{itemize}
  \item \pierre{The main setup does not isolate how much of the utility comes from
        the extra pretrained representation rather than the learnable memory.}
        Appendix D.2 uses an encoder trained on a disjoint 20\% split, but the
        main experiments use the ImageNet-pretrained one.
  \item \pierre{Compared against Retrain (M) only. It is not shown that deleting
        entries from a jointly trained memory equals training the whole
        architecture without them.}
  \item \pierre{Backbone ViT-Tiny ($\sim$5M parameters), key encoder ViT-Base
        ($\sim$86M). Fairer: a ViT-Tiny baseline that also gets the frozen
        ViT-Base features, without the learnable memory.}
\end{itemize}
\end{frame}

\begin{frame}{MUNKEY: criticism (2/3)}
\begin{itemize}
  \item \pierre{Missing simple baselines: frozen encoder plus a trained
        classifier; a ViT augmented with retrieved neighbour \emph{labels}
        instead of jointly learned tokens.}
  \item \pierre{No explicit retraining-equivalence or forgetting objective:
        $\theta$ likely still holds sample knowledge and may just be less
        exposed to this MIA.}
  \item \pierre{The average gap (over test, retain, forget accuracy and MIA AUC)
        can hide failures on a single metric.}
  \item \pierre{No distribution-shift tests, as needed for harmful data such as
        spam.}
\end{itemize}
\end{frame}

\begin{frame}{MUNKEY: criticism (3/3)}
\begin{itemize}
  \item \pierre{Forgetting is likely easier for \emph{redundant} examples: a
        neighbour then represents the forget sample well, with no visible
        forgetting cost.}
  \item \pierre{\textbf{Indirect influence through neighbour tokens.} Gradients
        into retrieved neighbour values are stopped, but the shared parameters
        $\theta$ are still updated. A sample $x_f$ influences $\theta$ through
        its own loss \emph{and} whenever its token is retrieved to train other
        samples. Deleting its entry undoes neither.}
\end{itemize}
\end{frame}

\begin{frame}{MUNKEY: what transfers to us, and what does not}
\begin{itemize}
  \item Our PKM is \emph{content}-addressed $\Rightarrow$ shared slots, no
        exclusivity. MUNKEY's guarantee comes from \emph{instance} addressing.
  \item Spam is a \emph{population} pattern learned by the backbone; deleting
        instance entries alone would not remove it.
  \item \pierre{We keep the items: an entry is \emph{updated} by unlearning, with
        emphasis on it over the backbone, not deleted.}
\end{itemize}
\end{frame}

\begin{frame}{Extensions implemented (1/4): memory dropout}
\textbf{Memory dropout on the PKM.} Each sequence's whole PKM read is replaced by
a learned null vector with probability $p$ during training, so the backbone works
without a user's memory contribution.
\begin{center}
\small
\begin{tabular}{|l|c|c|c|}
\hline
run & retain slots (subset) & Jaccard & test N@10 \\
\hline
no dropout       & 92k  & 0.43 & 0.0305 \\
dropout 0.1      & 106k & 0.41 & 0.0318 \\
dropout 0.3      & 101k & 0.44 & 0.0303 \\
\hline
\end{tabular}
\end{center}
Memory stays healthy, utility unchanged. \strun: unlearning comparison.
\end{frame}

\begin{frame}{Extensions implemented (2/4): item-keyed exemplar memory}
\textbf{Item-keyed exemplar memory} (replacing the PKM, or next to it). One
entry $v_i$ per item, added to the item's input tokens. Training per occurrence:
null (prob.\ 0.1), similarity-weighted average of 1--4 \textbf{neighbour items}
by frozen text embedding (prob.\ 0.2, no gradient), else own entry.
\end{frame}

\begin{frame}{Extensions implemented (3/4): training side}
\begin{enumerate}
  \item \textbf{Output-side item score} $s_i = \langle h_0, W v_i\rangle$
        ($h_0$ = decoder state after BOS), trained with a sampled softmax,
        used to rerank the beams. Forgetting can then push down the
        \emph{target's own} entry.
  \item \textbf{Item-addressed PKM slots:} every item owns a private block of 16
        slots next to the shared table $\Rightarrow$ exclusive by construction.
  \item \textbf{Consistency distillation}
        KL$(\text{sg}(p_{\text{own}})\,\|\,p_{\text{neighbours}})$: trains the
        backbone so neighbour substitution stays a good prediction.
\end{enumerate}
\end{frame}

\begin{frame}{Extensions implemented (4/4): unlearning side}
\begin{enumerate}
  \setcounter{enumi}{3}
  \item \textbf{Unlearning knobs:} lr emphasis (item entries $\times$10,
        backbone $\times$0.1 or frozen), rows = forget items / AF-IHF top-$t$ /
        soft AF-IHF weights, warm start towards the neighbour average.
  \item \textbf{Local repair as an additional objective term:} CE on retain rows
        whose history contains an affected item, \emph{next to} the unchanged
        uniform retain term.
\end{enumerate}
\end{frame}

% ===========================================================================
\section{Meeting notes and action items}

\begin{frame}{Meeting 31.07.2026}
\begin{itemize}
  \item Continual-learning critique: replay with old data is bad; rehearsal is
        a \textbf{privacy concern}.
  \item Forget-set-only unlearning: use statistics such as access counts;
        implement top-$t$ selection.
  \item Think about the connection between LSTMs and unlearning (gates).
  \item Document \emph{why} we insert memory and \emph{where}; use it as a
        modular component and show the difference.
  \item The Google doc describes how memory rows are selected in the
        continual-learning paper (next slides).
\end{itemize}
\end{frame}

\begin{frame}{Meeting 31.07.2026: status}
\begin{center}
\small
\begin{tabular}{|l|c|}
\hline
item & status \\
\hline
Top-$t$ selection from access statistics & \stdone \\
Why and where the memory goes (placement) & \stdone \\
Forget-set-only unlearning (no retain data)                          & \sttodo \\
LSTM gates and unlearning                                            & \sttodo \\
\hline
\end{tabular}
\end{center}

\medskip
\pierre{All current unlearning runs still use a retain term. Forget-only is
relevant for the privacy argument: no rehearsal of old data.}
\end{frame}

\begin{frame}{Memory-row selection in the continual-learning paper}
\textbf{Stage 1: learn $D_1$.} Update the backbone on the new data
$(q, \text{DocID}) \in D_1$ only; no $D_0$ examples are mixed in.

\medskip
\textbf{Stage 2: select memory rows.} Run the $D_1$ queries through the adapted
model and record which rows they access.
\begin{itemize}
  \item \textbf{Protect} rows used heavily by $D_0$: the top 10\% most
        historically used rows cannot change.
  \item Among the rest, \textbf{prefer} rows the $D_1$ queries access often and
        that were rarely used before (TF-IDF style).
  \item Select 10{,}000 rows.
\end{itemize}

\end{frame}

\begin{frame}{Memory-row selection in the continual-learning paper (cont.)}
\textbf{Stage 3: update only those rows.} Freeze the backbone, LoRA, keys,
routing, output embeddings and all unselected rows; train only the selected
value vectors on the correct $D_1$ DocID tokens. Then add the $D_1$ access
counts to the history and repeat for $D_2$.
\end{frame}

\begin{frame}{Their selection vs.\ ours}
\begin{center}
\small
\begin{tabular}{|l|l|l|}
\hline
 & continual learning & our unlearning \\
\hline
new data      & $D_1$ queries          & forget set \\
history       & $D_0$ access counts    & retain-set access counts \\
score         & now frequent, rarely used before & AF-IHF (same idea) \\
protection    & \textbf{top 10\% rows frozen} & none \\
rows          & 10{,}000               & $t \in \{100, 1000, 10^4\}$ \\
updated       & selected values only   & selected values only \\
\hline
\end{tabular}
\end{center}

\medskip
\begin{itemize}
  \item Selected rows $=$ useful for the new data $+$ unimportant for old data.
  \item \pierre{Not yet tried: the hard protection of the most-used retain
        rows. It directly targets the 4--6\% utility cost of memory-only
        unlearning.}
\end{itemize}
\end{frame}

\begin{frame}{Meeting 18.08.2026 (Kidist / Pierre)}
\begin{itemize}
  \item PKM is still useful as a fast lookup-based module, but given the FFN
        control we should \textbf{not overclaim} PKM as the mechanism.
  \item Safer story: \emph{localized decoder-layer rebuilding} removes the spam
        (freeze the model, reset only the first two decoder FFNs, retrain them
        on clean retain data); PKM gives a faster sparse implementation,
        especially under small update budgets.
  \item Next question: are the speed gains over a fresh FFN consistent across
        seeds, attacks and datasets?
\end{itemize}
\end{frame}

\begin{frame}{Meeting 18.08.2026 (Kidist / Pierre, cont.)}
\begin{itemize}
  \item With that story, top-$t$ selection is not the core method; only the
        placement result (d01) is needed. Selection matters only for a
        PKM-specific sparse-memory claim.
  \item Reference: \url{https://belinkov.com/assets/pdf/sigir2025.pdf}
\end{itemize}
\end{frame}

\begin{frame}{Meeting 18.08.2026: action items}
\begin{center}
\small
\begin{tabular}{|l|c|}
\hline
action item & status \\
\hline
Finish the FFN-reinit control                       & \stdone \\
Fix memory collapse                                 & \stdone \\
Top-$t$ sweeps on a healthy memory                  & \stdone\,$^{a}$ \\
AF-IHF vs.\ grad vs.\ grad+dot vs.\ access frequency & partly\,$^{b}$ \\
Adam vs.\ SGD under selection                       & \sttodo \\
Multi-seed experiments                              & partly\,$^{c}$ \\
Other attacks and datasets                          & partly\,$^{d}$ \\
\hline
\end{tabular}
\end{center}
{\footnotesize
$^{a}$ on the healthy \emph{trained-in} memory, not post-hoc.
$^{b}$ AF-IHF vs.\ grad+dot done; AF and grad alone not on the healthy memory.
$^{c}$ 3 seeds for the trained-in comparison; PKM-vs-FFN speed not yet.
$^{d}$ collapse fix on sports; unlearning only on beauty bandwagon.}
\end{frame}

\begin{frame}{Meeting: baselines and datasets}
\textbf{Trained from the start:}
\begin{itemize}
  \item ``Normal'' TIGER $+$ unlearning algorithms \hfill \strun
  \item Ours: TIGER with PKMs replacing some FFNs during training $+$
        unlearning \hfill \stdone
  \item PKM models must be competitive with normal training \hfill partly
        \pierre{(clean: 0.0313 balanced vs.\ 0.0321 without the fix)}
\end{itemize}

\end{frame}

\begin{frame}{Meeting: baselines and datasets (cont.)}
\textbf{PKM as adapter:}
\begin{itemize}
  \item ``Normal'' poisoned TIGER, unchanged \hfill \stdone
  \item TIGER with some FFNs re-initialised and retrained, rest frozen \hfill \stdone
  \item Ours: FFNs replaced by PKMs, PKMs trained, rest frozen \hfill \stdone
\end{itemize}

\medskip
\textbf{Datasets:} Amazon Beauty, Sports \& Outdoors, Toys \& Games.
\end{frame}

% ===========================================================================
\section{Plan (updated Oct 16)}

\begin{frame}{Status of the running work}
\begin{center}
\small
\begin{tabular}{|l|c|}
\hline
item & status \\
\hline
Routing balance loss (switch + batchnorm), replicated & \stdone \\
Slot-selection unlearning, budget sweep, full-model reference & \stdone \\
Seeds 3/4, 50 requests, sensitive $10^{-4}$ & \stdone \\
Memory dropout: training & \stdone \\
Memory dropout: unlearning comparison & \stdone \\
Item memory: replacing the PKM / next to the PKM & \stdone \\
Item memory + output score (+ consistency) & \stdone \\
Item-addressed PKM slots & \stdone \\
Unlearning variants: emphasis, soft rows, warm start & \stdone \\
Local repair (fixed: rows from the full retain split) & \strun \\
\hline
\end{tabular}
\end{center}
\end{frame}

\begin{frame}{Exploration plan (Oct 16): switch loss and MUNKEY}
\begin{center}
\small
\begin{tabular}{|l|l|c|}
\hline
\# & item & status \\
\hline
M1 & user tokens + history-pathway dropout ($p_{hist}$) & \stdone \\
M3 & whole-user deletion, 2\% of users (MUNKEY's setting) & \stdrop \\
M2 & pathway-sensitivity $P_s$ model selection & \stdrop \\
M4 & inference-time neighbour ensemble & \stdrop \\
S1 & anneal the switch loss to 0 (from step 1k / 2k) & \stdone \\
S5 & competitiveness: clean PKM vs.\ clean TIGER & \stdone \\
S2 & raw vs.\ standardised $P$; per-half vs.\ slot balancing & \stdone \\
S3 & placements: D01, E23D01, all layers & \stdone \\
S4 & capacity: $n_{keys}$ 256 / 1024 & \stdone \\
\hline
\end{tabular}
\end{center}
\end{frame}

\begin{frame}{Exploration plan: the two key tests (outcome)}
\begin{itemize}
  \item \textbf{M3, whole-user deletion:} delete 2\% of users. Their tokens are
        reset to the average of their nearest \emph{remaining} users, with zero
        gradient steps. Compared with a model retrained without them (served the
        same way) and with full-model unlearning, on deleted vs.\ other users.
        \alert{Dropped (Oct 19): user tokens cost 40--45\% utility at every setting.}
  \item \textbf{S1, annealing:} if the memory stays healthy after the loss is
        switched off, the switch loss is an initialisation aid, not a permanent
        component. \alert{Outcome: the router collapses again (128--144 slots), so
        the loss is needed throughout.}
\end{itemize}
\end{frame}

\begin{frame}{Results (Oct 19): switch-loss exploration}
Retain slots at eval (retain subset), E1 unless noted; reference: 92k slots,
N@10 0.0305.
\begin{center}
\small
\begin{tabular}{|l|c|c|c|}
\hline
run & retain slots & Jaccard & test N@10 \\
\hline
anneal from 1k / 2k & \alert{128 / 144} & 1.00 & .0309 / .0297 \\
raw $P$ & \alert{128}       & 1.00 & .0304 \\
$n_{keys}=256$ (65k slots)    & 26k (40\%)        & 0.42 & \textbf{.0328} \\
$n_{keys}=1024$ (1M slots)    & 81k (8\%)         & 0.62 & .0313 \\
D01, coef 0.3                 & 20k / 35k          & 0.57 & .0296 \\
E23D01 / ALL & 6--100k & 0.40--0.83 & .0316 / .0313 \\
\hline
\end{tabular}
\end{center}
\end{frame}

\begin{frame}{Results (Oct 19): what the switch exploration shows}
\begin{itemize}
  \item The loss is needed \textbf{throughout} training: annealing it away lets
        the router collapse again.
  \item Per-token \textbf{standardisation is essential}: raw $P$ collapses, as
        the scale-loophole argument predicts.
  \item A \textbf{smaller} memory is used better and gives the best utility.
  \item Multi-layer layouts are healthy, but far below single-layer E1.
\end{itemize}
\end{frame}

\begin{frame}{Results (Oct 19): item memory vs.\ user tokens}
Poisoned beauty, test N@10 / SH@10 (\alert{corrected Oct 23}: zero-step evals):
\begin{center}
\small
\begin{tabular}{|l|c|c|}
\hline
model & N@10 & SH@10 \\
\hline
plain TIGER                         & .0315 & .0126 \\
item memory                         & .0352 & .0120 \\
\quad + output score               & .0364 & .0200 \\
\quad + score + consistency        & \textbf{.0382} & .0120 \\
\quad + score + consistency + PKM  & .0375 & .0211 \\
\hline
user tokens (any $p_{hist}$, $p_{drop}$) & \alert{.017--.020} & --- \\
\hline
\end{tabular}
\end{center}
\begin{itemize}
  \item Item memory: \textbf{+12 to +21\% utility}. \alert{Withdrawn (Oct 23):}
        ``less spam pickup'' came from mis-evaluated before rows.
  \item User tokens: $-$40 to $-$45\% \pierre{(likely memorisation of past
        targets)}. \textbf{Decision (Oct 19): user tokens dropped.}
\end{itemize}
\end{frame}

\begin{frame}{Results (Oct 19): memory dropout has no effect}
Unlearning on the balanced E1 model, best memory-only arm vs.\ full model (N@10):
\begin{center}
\small
\begin{tabular}{|l|c|c|c|}
\hline
memory dropout & full model & best memory-only & gap \\
\hline
none   & .0313 & .0294 & $-$6\% \\
$p=0.1$ & .0330 & .0308 & $-$7\% \\
$p=0.3$ & .0309 & .0297 & $-$4\% \\
\hline
\end{tabular}
\end{center}
The gap to full-model unlearning stays (1 seed), and at $p=0.3$ removal drops.
\end{frame}

\begin{frame}{Results (Oct 20): item-memory unlearning}
Old recipe ($\lambda_f$ 0.01, $\lambda_s$ 1.0), test N@10 / SH@10; \emph{emph} =
entries $\times$10, backbone $\times$0.1; before column \alert{corrected Oct 23}:
\begin{center}
\footnotesize\setlength{\tabcolsep}{3pt}
\begin{tabular}{|l|c|c|c|c|}
\hline
model & before & full model & emph & entries only \\
\hline
item memory           & .0352/.0120 & .0354/0 & .0353/.0003 & .0353/.0115 \\
\quad + score        & .0364/.0200 & .0376/0 & .0365/.0001 & .0364/.0199 \\
\quad + score + cons. & .0382/.0120 & .0377/0 & \textbf{.0382}/.0001 & .0382/.0120 \\
\quad + PKM          & .0375/.0211 & .0381/0 & .0373/.0002 & .0375/.0207 \\
\hline
\end{tabular}
\end{center}
\end{frame}

\begin{frame}{Results (Oct 20): item-memory unlearning, takeaways}
\begin{itemize}
  \item \textbf{Emphasis works with this recipe:} 97--100\% removed at
        full-model utility. \alert{Oct 23: only because of the strong separation
        term ($\lambda_s = 1$); with the paper recipes it is slow (see below).}
  \item \textbf{Entries alone do not:} with the backbone frozen nothing is
        removed. The small backbone step is required.
  \item Soft rows = hard rows; the warm-start reset costs 1--3\%.
  \item Local repair: fixed Oct 21 (rows from the full retain split); no
        effect, because unlearning already sits above the retrained utility.
\end{itemize}
\end{frame}

\begin{frame}{Decisions (Oct 20)}
\begin{itemize}
  \item Item memory with emphasis is the main unlearning recipe.
        \alert{(Revisited Oct 23: the gain is the item memory, not emphasis.)}
  \item Switch loss (0.1) + batchnorm query are now the PKM training default.
  \item Soft row selection preferred (modelling view; same numbers as hard rows).
  \item Dropped: warm-start reset (costs utility), top-10\% slot protection
        (blocks removal), user tokens.
\end{itemize}
\end{frame}

\begin{frame}{Local repair fix (Oct 20)}
\begin{itemize}
  \item Local rows are now drawn from the \emph{full} retain split: up to 256
        rows whose history has an affected item (seeded), one extra batch per
        step, same retain CE, weight $\lambda_{local}$.
  \item The uniform retain term is unchanged; only the 256 rows are forwarded.
  \item Reruns on all five item-memory models queued.
\end{itemize}
\end{frame}

\begin{frame}{Results (Oct 20): meeting action items}
Balanced E1 model, before: .0305 / .0052 (N@10 / SH@10):
\begin{center}
\footnotesize\setlength{\tabcolsep}{3pt}
\begin{tabular}{|l|c|c|}
\hline
arm & N@10 & SH@10 \\
\hline
AF-IHF $t$=1000 (reference)        & .0294 & .0005 \\
\quad + protect top 10\% slots   & .0303 & \alert{.0061} \\
forget-only, AF-IHF $t$=1000       & .0287 & .0002 \\
forget-only, full model            & \alert{.0000} & 0 \\
SGD lr 0.1 / 1.0 (AF-IHF)          & .0309 / .0299 & \alert{.0146 / .0051} \\
\hline
\end{tabular}
\end{center}
\begin{itemize}
  \item Protection blocks removal: the spam sits in the most retain-used slots.
  \item \textbf{Without retain data, memory-only still works; the full model is
        destroyed} (a privacy argument). Adam beats SGD here.
\end{itemize}
\end{frame}

\begin{frame}{MUNKEY: what the authors' code actually does}
\begin{itemize}
  \item One 128-d token per \emph{training instance} (+ linear adapter, null
        token); frozen keys from a pretrained ViT-Base, used for retrieval only.
  \item Training, decided \emph{once per batch}: w.p.\ $p_r$ the token is the
        softmax$(\mathrm{sim}/0.07)$ mix of $K' \sim U\{2..k\}$ neighbours (the
        whole batch excluded, neighbours detached); w.p.\ $\alpha$ null token,
        w.p.\ $\beta$ the image is masked.
  \item Inference \emph{always} retrieves (top-4); the README's main command
        mixes tokens in one pass, not the paper's $K$-pass ensemble.
  \item Unlearning: mask the forget keys at retrieval (zero-shot).
  \item Code defaults switch all of this off; the paper uses $p_r = 0.2$ and
        picks $(\alpha, \beta)$ by $P_s$.
\end{itemize}
\end{frame}

\begin{frame}{MUNKEY: baseline port and recipe transfer (results Oct 23)}
\begin{enumerate}
  \item \textbf{Baseline (faithful):} instances = training users, keys = frozen
        text-embedding mean of the history, token prepended to the encoder.
        Eval mixes the nearest \emph{other} users (never the own token: the
        failure mode of our user tokens). Unlearning = delete the fake users'
        keys. Grid $\alpha = \beta \in \{0, 0.1, 0.2\}$, plus a clean
        retrained reference.
  \item \textbf{Our base:} item memory trained with MUNKEY's exact recipe
        (per-batch draws, $\tau = 0.07$, eval retrieval over [self +
        neighbours]); then our emphasis unlearning.
  \item Arms per model: key deletion only (zero-shot), key deletion + full
        unlearning, full unlearning alone.
\end{enumerate}
\end{frame}

\begin{frame}{Planned next (updated Oct 20)}
\begin{enumerate}
  \item Item-memory unlearning replicated: \stdone{} (data seeds 3, 5).
  \item Local repair with rows from the full retain split: \stdone{} (no effect).
  \item $n_{keys}=256$: \stdone{} --- does not replicate, 512 stays.
  \item Sensitive $10^{-3}$: \stdrop{} for now (spam only, Oct 21).
  \item Sports: \stdone{} --- the spam attack did not land there.
\end{enumerate}
\end{frame}

\begin{frame}{Queued (Oct 19/20): status}
\begin{center}
\footnotesize\setlength{\tabcolsep}{4pt}
\begin{tabular}{|l|l|c|}
\hline
goal & runs & status \\
\hline
replicate item memory & seeds 3, 4: full / emph & \stdone \\
replicate $n_{keys}=256$ & seeds 3, 4 (+ slot diagnostics) & \stdone \\
item memory, sens.\ $10^{-3}$ & before / full / emph, retrain ref. & \stdrop \\
second dataset & sports: plain vs.\ item memory & \stdone \\
local repair, fixed & emph + local, 5 models & \stdone \\
protect top 10\% retain slots & AF-IHF $t$=1000 / 100 & \stdone \\
forget-only unlearning & values, AF-IHF, full model & \stdone \\
Adam vs.\ SGD under selection & SGD lr 0.1 / 1.0 & \stdone \\
\hline
\end{tabular}
\end{center}
\end{frame}

\begin{frame}{Planned next (from the meeting notes)}
\begin{enumerate}
  \setcounter{enumi}{5}
  \item Protect the most-used retain rows (top 10\%) during selection, as in the
        continual-learning paper. \stdone
  \item Forget-set-only unlearning (no retain data, no rehearsal). \stdone
  \item Adam vs.\ SGD under selection. \stdone
  \item PKM vs.\ fresh FFN speed gain across seeds, attacks and datasets.
  \item ``Normal'' TIGER $+$ unlearning baseline: \stdone{} (poisoned .0313, after full unlearning .0325).
\end{enumerate}

\medskip
\pierre{Open question for the paper: is the contribution the collapse fix, the
locality of memory-only unlearning, or the item-keyed memory?}
\end{frame}

\section{Results (Oct 21 to 23)}

\begin{frame}{Paper recipes and references}
\begin{itemize}
  \item \textbf{Forget + Repair (F+R):} $\lambda_f = 0.1$, $\lambda_s =
        \lambda_n = 0$ (the paper's validation-selected setting).
  \item \textbf{NAU:} $\lambda_f = 0.1$, $\lambda_s = 0.01$, $\lambda_n =
        0.01$, 8 text-embedding neighbors, mass loss.
  \item Both match the recorded paper runs setting by setting; the NAU code on
        this branch matches main (audited; validation-based epoch selection
        ported, off by default).
  \item \textbf{Retrained references} (clean-trained models, scored against the
        spam targets): plain .0305 / 0, item memory .0363 / 0, balanced PKM
        .0301 / 0.
\end{itemize}
\end{frame}

\begin{frame}{Full unlearning with the paper recipes}
Test N@10 / SH@10; seeds are the paper's data seeds (own target, own fake users):
\begin{center}
\footnotesize\setlength{\tabcolsep}{4pt}
\begin{tabular}{|l|c|c|c|}
\hline
model, seed & before & full F+R & full NAU \\
\hline
plain, 2       & .0315/.0126 & .0328/.0003 & .0326/.0002 \\
item mem., 2   & .0382/.0120 & .0385/.0003 & .0385/0 \\
plain, 3       & .0323/.0012 & .0331/0 & .0325/0 \\
item mem., 3   & .0377/.0032 & .0382/0 & .0381/0 \\
plain, 5       & .0341/.0201 & .0348/0 & .0344/0 \\
item mem., 5   & .0382/.0112 & .0377/0 & .0384/0 \\
\hline
plain, clone-inject     & .0318/.0033 & .0323/.0003 & .0320/0 \\
item mem., clone-inject & .0381/.0017 & .0378/0 & .0380/0 \\
\hline
\end{tabular}
\end{center}
\end{frame}

\begin{frame}{Full unlearning: takeaways}
\begin{itemize}
  \item Both paper recipes remove the spam on every model.
  \item \textbf{Item memory keeps an 8 to 17\% utility lead} over plain after
        unlearning, at all three data seeds and on the clone-inject attack.
  \item Unlearned models sit above their retrained utility (item memory .0385
        vs .0363), since unlearning includes a retain fine-tune.
  \item NAU and F+R are close on Beauty
  \item The seed-3 attack barely landed (plain SH .0012): little to remove.
\end{itemize}
\end{frame}

\begin{frame}{Emphasis and zero-shot deletion}
\begin{itemize}
  \item \textbf{Emphasis} (entries $\times$10, backbone $\times$0.1) with the
        paper recipes: 25 to 40\% removal at 4 epochs; validation-selected
        epochs climb to the 10-epoch cap (about 70\%). Slow, not broken.
  \item \textbf{The utility gain is the architecture:} full F+R on item memory
        (.0385) matches the best emphasis result.
  \item \textbf{Zero-shot item deletion} (target's entry becomes its neighbors'
        mix): SH unchanged. The spam lives in the backbone, not in the target's
        entry. \pierre{This confirms kind of that the MUNKEY results stem from the easy datasets and that the params also contain knowledge that we might need to touch}
  \item \textbf{Local repair:} no effect; there is no utility gap left to repair.
\end{itemize}
\end{frame}

\begin{frame}{MUNKEY baseline (faithful port)}
\begin{itemize}
  \item \textbf{Spam:} key deletion of the fake users removes only 8 to 20\%
        (SH .0174 to .0140 at the best setting); key deletion plus full
        unlearning removes everything. MUNKEY models have about 25\% lower
        utility than plain (N@10 .023 to .025).
  \item \textbf{Whole-user deletion} (its home setting, random 2\% of users):
        membership-inference AUC .699 before, .613 after key deletion, .564
        retrained: 64\% of the gap closed with no weight update, retained users
        unaffected (N@10 .0247 to .0245).
  \item Deleted users' own N@10 rises after deletion (.0289 to .0306): their own
        token had memorized training targets, as with our user tokens.
  \item MUNKEY's training recipe on our item memory costs 5 to 7\% utility.
\end{itemize}
\end{frame}

\begin{frame}{Switch loss: impact on utility}
Test N@10, collapsed vs switch loss (loss term coefficient 0.1 and use batchnorm):
\begin{center}
\footnotesize\setlength{\tabcolsep}{4pt}
\begin{tabular}{|l|c|c|c|}
\hline
setting & collapsed & switch & change \\
\hline
Beauty, encoder block 1, seed 2 & .0307 & .0305 & $-$0.7\% \\
training seed 3                 & .0318 & .0320 & +0.6\% \\
training seed 4                 & .0326 & .0328 & +0.6\% \\
clean data                      & .0321 & .0313 & $-$2.5\% \\
decoder blocks 0, 1 (1 seed)    & .0314 & .0299 & $-$4.8\% \\
Sports, encoder block 1 (1 seed) & .0209 & .0188 & $-$10\% \\
\hline
\end{tabular}
\end{center}
Neutral in the main setting (inside seed noise); the decoder and Sports costs
are single runs. The memory reads 92k to 165k of 262k slots instead of 32.
\end{frame}

\begin{frame}{Switch loss: all positions and smaller results}
\begin{itemize}
  \item \textbf{All 8 FFNs as PKMs:} utility fine (.0313), but the loss is the
        \emph{mean} over memories, so each layer got 1/8 of the pressure:
        encoder memories read 6 to 8k slots, forget/retain overlap .80.
        Matched pressure (coef 0.8) is training.
  \item $n_{keys}=256$ does not replicate (mean +1.5\%); 512 stays.
  \item \textbf{50 requests:} full unlearning loses 3.6 to 4.7\% utility.
  \item \textbf{Validation-selected epochs:} full unlearning picks the fixed 4.
  \item \textbf{Cost:} every method takes the same 40 steps; memory methods buy
        locality (parameters allowed to move), not speed.
  \item Mean-output ablation of collapsed vs balanced memory: running.
\end{itemize}
\end{frame}

\section{Results (Oct 26): paper grid}

\begin{frame}{Paper grid: item memory vs plain (15 cells)}
3 target strata $\times$ data seeds 2, 3, 5, 7, 11; test N@10, plain TIGER plus an explicit per-item memory (``im'') with an output-score rerank (``s'') and a consistency loss (``c'') (short: paired imsc) vs plain:
\begin{center}
\footnotesize\setlength{\tabcolsep}{4pt}
\begin{tabular}{|l|c|c|c|c|}
\hline
 & plain & imsc & imsc vs plain & min / max \\
\hline
before deletion   & .0313 & .0379 & +21.4\% & +7.9 / +33.6\% \\
after full F+R    & .0319 & .0381 & \textbf{+19.5\%} & +8.1 / +28.9\% \\
after full NAU    & .0319 & .0381 & \textbf{+19.8\%} & +8.0 / +28.6\% \\
\hline
\end{tabular}
\end{center}
\begin{itemize}
  \item Positive in every cell; both recipes remove the spam on both
        architectures (SH $\leq$ .0004 almost everywhere).
  \item Popular targets: the retrained model still recommends them (reference SH
        up to .013), so unlearning to SH $\approx$ 0 over-suppresses, for both
        architectures alike (an objective issue, not an architecture one).
\end{itemize}
\end{frame}

\begin{frame}{Collapsed vs balanced PKM under the paper recipes}
Encoder block 1, test N@10 / SH@10 (fresh zero-step before evals):
\begin{center}
\footnotesize\setlength{\tabcolsep}{3pt}
\begin{tabular}{|l|c|c|c|}
\hline
model & before & full F+R / NAU & memory-only F+R / NAU \\
\hline
collapsed, s2 & .0293/.0091 & .0304/0, .0304/0 & .0297/.0057, .0297/.0065 \\
balanced, s2  & .0309/.0161 & .0313/.0001, .0312/0 & .0306/.0116, .0305/.0127 \\
collapsed, s3 & .0315/.0114 & .0321/.0001, .0317/0 & .0313/.0084, .0314/.0090 \\
balanced, s3  & .0310/.0086 & .0327/0, .0320/0 & .0313/.0076, .0311/.0078 \\
\hline
\end{tabular}
\end{center}
\begin{itemize}
  \item Full unlearning works on both; memory-only removes only 10 to 35\% on
        both. The balanced memory gives no unlearning advantage here.
  \item The earlier memory-only success needed $\lambda_s = 1$, like emphasis.
\end{itemize}
\end{frame}

\begin{frame}{All positions with matched pressure}
Switch coefficient 0.8 over 8 memories (= 0.1 per layer, as for one layer):
\begin{itemize}
  \item Encoder memories now read 18 to 29\% of their slots (2 to 3\% at
        coefficient 0.1), forget/retain overlap .35 to .66 (was .80), 13k to 16k
        forget-exclusive slots per layer (was 150 to 394).
  \item Decoder memories read 9 to 16\%.
  \item Utility comparable to one layer (training-time test N@10 .0315).
\end{itemize}
\medskip
\pierre{Collapse is fixable everywhere; the open question is whether a
separable memory can be exploited by an unlearning algorithm at all.}
\end{frame}

\begin{frame}{Does the spam live in the memory? Slot deletion (Oct 26)}
Delete slots at evaluation (never selected in top-k); test N@10 / SH change:
\begin{center}
\footnotesize\setlength{\tabcolsep}{3pt}
\begin{tabular}{|l|c|c|c|}
\hline
slots deleted per memory & encoder block 1 & all 8 positions & forget reads \\
\hline
strictly forget-exclusive (60--65) & 0 / 0 & 0 / 0 & $\approx$0\% \\
top-1000 forget-specific & $-$0.2\% / $-$9\% & +0.2\% / $-$2\% & 49\% \\
\textbf{every slot the fake users read} & $-$0.4\% / $-$13\% & +1.4\% / $-$2\% & 100\% \\
\hline
\end{tabular}
\end{center}
\begin{itemize}
  \item Strict separability vanishes against the full retain set.
  \item Deleting every forget-read slot removes only 2 to 13\% of the spam: it is
        stored in attention and the backbone, not in the memory. A routing loss
        cannot exploit separability for this spam.
\end{itemize}
\end{frame}

\begin{frame}{Boosted item entries with an anti-nuking cap (Oct 26)}
Backbone 1$\times$, forget-read entries 10$\times$; each entry capped at $c \cdot
\lVert v_0 \rVert$ from its pre-unlearning value over the whole sequence:
\begin{center}
\footnotesize\setlength{\tabcolsep}{3pt}
\begin{tabular}{|l|c|c|c|c|}
\hline
cell (retrained SH) & full & boost & cap 0.5 & cap 0.2 \\
\hline
mid s2 (0)          & .0385/.0003 & .0385/.0002 & .0385/.0002 & .0385/.0002 \\
popular s3 (.0097)  & .0376/0 & .0377/0 & .0377/0 & .0377/0 \\
popular s11 (.0130) & .0390/0 & .0394/0 & .0392/0 & .0394/0 \\
\hline
\end{tabular}
\end{center}
\begin{itemize}
  \item Harmless to slightly positive; entries drift at most 0.4 of their norm,
        so nothing is nuked and the cap is a free guarantee.
  \item Popular targets end below the retrained SH either way: the backbone's
        forget term over-suppresses, not the entries.
\end{itemize}
\end{frame}

\begin{frame}{Forget-dominated, not exclusive, slots (Oct 26)}
Slots the fake users read at $\geq r\times$ the per-row retain rate, deleted at
evaluation; test N@10 / SH (share of forget reads / retain reads deleted):
\begin{center}
\footnotesize\setlength{\tabcolsep}{3pt}
\begin{tabular}{|l|c|c|}
\hline
deleted & encoder block 1 & all 8 positions \\
\hline
none & .0275 / .0163 & .0296 / .0073 \\
$r=10$ & .0275 / .0163 \ (3\% / 0.2\%) & .0296 / .0073 \ (1\% / 0.1\%) \\
$r=1$ & .0276 / .0150 \ (78\% / 40\%) & .0296 / .0075 \ (70\% / 45\%) \\
\hline
\end{tabular}
\end{center}
\begin{itemize}
  \item Even deleting the slots that carry three quarters of the fake users'
        reads changes SH by $-$8\% / +3\%: dominance does not locate the spam either.
  \item Utility does not move: deleted slots are replaced by the next top-$k$.
  \item Deletion is not the goal anyway (we keep recommending the item);
        next: update dominated slots \emph{differently} (next slide).
\end{itemize}
\end{frame}

\begin{frame}{Running (Oct 26)}
\begin{itemize}
  \item \textbf{Dominated-slot boost} (PKM): full F+R / NAU, but the top-1000
        AF-IHF slots per memory move at 10$\times$, each capped at 0.5 of its
        norm (only boosted slots are capped). Encoder block 1 (seeds 2, 3),
        all 8 positions (seed 2) + its pre / full references.
  \item \textbf{FFN pathway dropout} (add mode, $p=0.5$ per sequence): pushes
        the layer's job into the memory; then rerun slot deletion. Does the
        spam move into the memory?
  \item \textbf{Category memory:} keyed by the first SID level (103 groups on
        Beauty), 32 or 128 slots per group, alone and with the item memory;
        pre + F+R / NAU.
  \item \textbf{Balance loss} now sums over memories (coefficient 0.1 per
        memory, scales with the number of PKM layers).
\end{itemize}
\end{frame}

\begin{frame}{Why a collapsed PKM keeps its utility (Oct 23)}
Memory output replaced by its mean or by zero at evaluation, test N@10 change:
\begin{center}
\footnotesize\setlength{\tabcolsep}{4pt}
\begin{tabular}{|l|c|c|c|}
\hline
model & output vs input & token variation & mean / zero \\
\hline
collapsed, seed 2 & 0 (exactly) & none & 0 / 0 \\
collapsed, seed 3 & 71{,}000$\times$ & none & 0 / +0.6\% \\
balanced, seed 2  & 0.21$\times$ & strong & $-$0.2 / +0.3\% \\
balanced, seed 3  & 0.17$\times$ & strong & $-$0.2 / $-$1.2\% \\
\hline
\end{tabular}
\end{center}
\begin{itemize}
  \item A collapsed memory outputs the same constant for every token: no token
        information; the residual connection carries the network around it.
  \item Even the balanced memory's rich output can be removed for $\leq$1.2\%:
        this FFN position matters little for Beauty's ranking.
\end{itemize}
\end{frame}

\begin{frame}{Decisions (Oct 23)}
\begin{itemize}
  \item \textbf{Framing:} this is the PKM paper; we change the architecture. F+R
        and NAU are the unlearning algorithms run on each architecture; MUNKEY is the architectural baseline.
  \item Headline: item memory; the collapse fix is a supporting result.
  \item Dropped: emphasis, zero-shot item deletion, local repair.
  \item MUNKEY setting chosen on validation; epoch selection stays off.
  \item \textbf{Paper grid running:} plain + item memory on 3 target strata
        $\times$ 5 data seeds (24 new trainings, references per cell).
\end{itemize}
\end{frame}

\begin{frame}{Decisions pending (Oct 23, now decided: see above)}
\begin{enumerate}
  \item \textbf{Emphasis:} drop it and present item memory + standard
        unlearning (F+R / NAU), or tune it (milder backbone damping, more epochs)?
  \item \textbf{Contribution:} the item-memory architecture (utility lead that
        survives unlearning), the collapse fix (separable PKM), or both?
  \item \textbf{MUNKEY baseline:} which pathway-dropout setting to report
        (0 / 0.1 / 0.2 differ in utility and removal)?
  \item \textbf{Paper grid:} extend plain + item memory to data seeds 7, 11
        and the unpopular / popular targets (about 16 trainings)?
\end{enumerate}
\end{frame}

\begin{frame}{Decisions pending (Oct 23, cont.)}
\begin{enumerate}
  \setcounter{enumi}{4}
  \item \textbf{Second dataset:} the Sports attack did not land at the mid
        target; try another target stratum there, or use Toys?
  \item \textbf{Drop} zero-shot item deletion and local repair (both inert)?
  \item \textbf{Validation-based epoch selection:} keep it off (it picks the
        fixed 4 epochs for full unlearning)?
  \item \textbf{Sensitive scenario:} when to resume it?
\end{enumerate}
\end{frame}

\end{document}



